% !TEX program = lualatex
% ---------------------------------------------------------------------------------------------
%  Deformalization 3: descent to a layer of the cyclotomic tower   (riemann-rev)
%
%  Build:  python3 -B scripts/build_digest.py layer-descent
%
%  One unit of the Lean archive behind L. Alpöge's announced proof of RH, rewritten as ordinary mathematics.
%  The prose is hand-written. Lean names, line ranges and counts come from data.tex; the quoted Lean statements
%  are copied from the archive by the script. The proofs were written from the reading project's structured
%  proofs; the formal statements and the file of definitions were read in the archive.
%
%  Rules of the text: standard terminology; every symbol is explained where it first appears and listed again
%  in the notation at the end; no letter has two meanings. The signature rule: a function that the archive
%  defines is written with every explicit argument of its Lean definition, through the macros that the build
%  writes into data.tex from digests/notation.json and accounts.json (this file defines none of them), and every
%  hypothesis stands in its statement or in a Setting that the statement cites. Letters: the archive's subgroup U
%  is H, its twist exponent e is h, its coefficient ring R is Lambda, its fields k and L are F and E.
%
%  Layout: the reading edition of the guide (11 pt Libertinus on a 140 mm measure); preamble as in the first note.
% ---------------------------------------------------------------------------------------------
\documentclass[11pt,a4paper]{article}
\ifdefined\pdfvariable\pdfvariable suppressoptionalinfo \numexpr 1+2+4+8\relax\fi
\usepackage[a4paper,hmargin=35mm,top=22mm,bottom=27mm,footskip=12mm,nohead]{geometry}

\usepackage{fontspec}
\directlua{
  luaotfload.add_fallback("gfallback", {"STIXTwoMath-Regular.otf:mode=node;", "DejaVuSans.ttf:mode=node;"})
  luaotfload.add_fallback("gmonofallback", {"DejaVuSansMono.ttf:mode=node;", "STIXTwoMath-Regular.otf:mode=node;"})
}
\setmainfont{LibertinusSerif}[Extension=.otf,UprightFont=*-Regular,ItalicFont=*-Italic,
  BoldFont=*-Semibold,BoldItalicFont=*-SemiboldItalic,RawFeature={fallback=gfallback}]
\setsansfont{Inter}[Extension=.otf,UprightFont=*-Regular,ItalicFont=*-Italic,
  BoldFont=*-SemiBold,BoldItalicFont=*-SemiBoldItalic,RawFeature={fallback=gfallback}]
\setmonofont{JetBrainsMono}[Extension=.otf,UprightFont=*-Regular,ItalicFont=*-Italic,
  BoldFont=*-SemiBold,BoldItalicFont=*-SemiBoldItalic,Scale=MatchLowercase,
  RawFeature={-calt;-liga;fallback=gmonofallback}]
\usepackage{amsmath}
\usepackage{amsthm}
\usepackage[math-style=TeX]{unicode-math}
\setmathfont{LibertinusMath-Regular.otf}

\usepackage{xcolor}
\usepackage{tikz}
\usetikzlibrary{calc,arrows.meta,positioning,backgrounds}
\usepackage[protrusion=true,expansion=true]{microtype}
\usepackage{enumitem}
\usepackage{needspace}
\usepackage{booktabs}
\usepackage{array}
\usepackage{longtable}
\usepackage{fancyhdr}
\usepackage{lastpage}
\usepackage{multicol}
\usepackage[hidelinks,pdfusetitle]{hyperref}

\definecolor{ink}{HTML}{0F172A}
\definecolor{inksoft}{HTML}{475569}
\definecolor{inkfaint}{HTML}{94A3B8}
\definecolor{rulegrey}{HTML}{CBD5E1}
\definecolor{arith}{HTML}{C2410C}
\definecolor{est}{HTML}{7C3AED}
\definecolor{ana}{HTML}{0891B2}
\definecolor{ok}{HTML}{047857}
\definecolor{leancol}{HTML}{3730A3}

\title{Descent to a layer of the cyclotomic tower}
\author{riemann-rev}
\date{}

% ------------------------------------------------------------------------------------- macros
\newcommand{\leanfont}{\ttfamily\fontsize{8.6}{10.5}\selectfont\color{leancol}}
\newcommand{\lean}[1]{\csname LN@#1\endcsname}
\newcommand{\leanquote}[1]{\csname LQ@#1\endcsname}
% generated by scripts/build_digest.py -- do not edit
\expandafter\newcommand\csname LN@dRep\endcsname{\href{https://github.com/Deicyde/riemann-hypothesis/blob/40e56bab026552e089d7afcfeff08d2defb9c005/src/Development/AutomorphicArithmetic/AutomorphicVarieties/NumberFieldTowers/XTwApplyApply.lean\#L41-L43}{\leanfont Tower.dRep}}
\expandafter\newcommand\csname LN@xRep\endcsname{\href{https://github.com/Deicyde/riemann-hypothesis/blob/40e56bab026552e089d7afcfeff08d2defb9c005/src/Development/AutomorphicArithmetic/AutomorphicVarieties/NumberFieldTowers/XTwApplyApply.lean\#L45-L46}{\leanfont Tower.xRep}}
\expandafter\newcommand\csname LN@xRepApply\endcsname{\href{https://github.com/Deicyde/riemann-hypothesis/blob/40e56bab026552e089d7afcfeff08d2defb9c005/src/Development/AutomorphicArithmetic/AutomorphicVarieties/NumberFieldTowers/XTwApplyApply.lean\#L48-L48}{\leanfont Tower.xRep\_apply}}
\expandafter\newcommand\csname LN@twist\endcsname{\href{https://github.com/Deicyde/riemann-hypothesis/blob/40e56bab026552e089d7afcfeff08d2defb9c005/src/Development/AutomorphicArithmetic/AutomorphicVarieties/NumberFieldTowers/XTwApplyApply.lean\#L50-L54}{\leanfont Tower.twist}}
\expandafter\newcommand\csname LN@dTw\endcsname{\href{https://github.com/Deicyde/riemann-hypothesis/blob/40e56bab026552e089d7afcfeff08d2defb9c005/src/Development/AutomorphicArithmetic/AutomorphicVarieties/NumberFieldTowers/XTwApplyApply.lean\#L56-L56}{\leanfont Tower.dTw}}
\expandafter\newcommand\csname LN@xTw\endcsname{\href{https://github.com/Deicyde/riemann-hypothesis/blob/40e56bab026552e089d7afcfeff08d2defb9c005/src/Development/AutomorphicArithmetic/AutomorphicVarieties/NumberFieldTowers/XTwApplyApply.lean\#L58-L58}{\leanfont Tower.xTw}}
\expandafter\newcommand\csname LN@dTwApply\endcsname{\href{https://github.com/Deicyde/riemann-hypothesis/blob/40e56bab026552e089d7afcfeff08d2defb9c005/src/Development/AutomorphicArithmetic/AutomorphicVarieties/NumberFieldTowers/XTwApplyApply.lean\#L60-L61}{\leanfont Tower.dTw\_apply}}
\expandafter\newcommand\csname LN@xTwApply\endcsname{\href{https://github.com/Deicyde/riemann-hypothesis/blob/40e56bab026552e089d7afcfeff08d2defb9c005/src/Development/AutomorphicArithmetic/AutomorphicVarieties/NumberFieldTowers/XTwApplyApply.lean\#L63-L64}{\leanfont Tower.xTw\_apply}}
\expandafter\newcommand\csname LN@xTwApplyApply\endcsname{\href{https://github.com/Deicyde/riemann-hypothesis/blob/40e56bab026552e089d7afcfeff08d2defb9c005/src/Development/AutomorphicArithmetic/AutomorphicVarieties/NumberFieldTowers/XTwApplyApply.lean\#L66-L69}{\leanfont Tower.xTw\_apply\_apply}}
\expandafter\newcommand\csname LN@ShaD\endcsname{\href{https://github.com/Deicyde/riemann-hypothesis/blob/40e56bab026552e089d7afcfeff08d2defb9c005/src/Development/AutomorphicArithmetic/AutomorphicVarieties/NumberFieldTowers/XTwApplyApply.lean\#L73-L74}{\leanfont Tower.ShaD}}
\expandafter\newcommand\csname LN@IsGen\endcsname{\href{https://github.com/Deicyde/riemann-hypothesis/blob/40e56bab026552e089d7afcfeff08d2defb9c005/src/Development/AutomorphicArithmetic/AutomorphicVarieties/NumberFieldTowers/XTwApplyApply.lean\#L36-L37}{\leanfont Tower.IsGen}}
\expandafter\newcommand\csname LN@cZero\endcsname{\href{https://github.com/Deicyde/riemann-hypothesis/blob/40e56bab026552e089d7afcfeff08d2defb9c005/src/Development/AutomorphicArithmetic/AutomorphicVarieties/NumberFieldTowers/XTwApplyApply.lean\#L104-L104}{\leanfont Tower.cZero}}
\expandafter\newcommand\csname LN@layU\endcsname{\href{https://github.com/Deicyde/riemann-hypothesis/blob/40e56bab026552e089d7afcfeff08d2defb9c005/src/Development/AutomorphicArithmetic/AutomorphicVarieties/NumberFieldTowers/XTwApplyApply.lean\#L108-L108}{\leanfont Tower.layU}}
\expandafter\newcommand\csname LN@layF\endcsname{\href{https://github.com/Deicyde/riemann-hypothesis/blob/40e56bab026552e089d7afcfeff08d2defb9c005/src/Development/AutomorphicArithmetic/AutomorphicVarieties/NumberFieldTowers/XTwApplyApply.lean\#L110-L110}{\leanfont Tower.layF}}
\expandafter\newcommand\csname LN@XN\endcsname{\href{https://github.com/Deicyde/riemann-hypothesis/blob/40e56bab026552e089d7afcfeff08d2defb9c005/src/Development/AutomorphicArithmetic/AutomorphicVarieties/NumberFieldTowers/XTwApplyApply.lean\#L112-L112}{\leanfont Tower.XN}}
\expandafter\newcommand\csname LN@XX\endcsname{\href{https://github.com/Deicyde/riemann-hypothesis/blob/40e56bab026552e089d7afcfeff08d2defb9c005/src/Development/AutomorphicArithmetic/AutomorphicVarieties/NumberFieldTowers/XTwApplyApply.lean\#L114-L114}{\leanfont Tower.XX}}
\expandafter\newcommand\csname LN@XP\endcsname{\href{https://github.com/Deicyde/riemann-hypothesis/blob/40e56bab026552e089d7afcfeff08d2defb9c005/src/Development/AutomorphicArithmetic/AutomorphicVarieties/NumberFieldTowers/XTwApplyApply.lean\#L116-L116}{\leanfont Tower.XP}}
\expandafter\newcommand\csname LN@GenLE\endcsname{\href{https://github.com/Deicyde/riemann-hypothesis/blob/40e56bab026552e089d7afcfeff08d2defb9c005/src/Development/AutomorphicArithmetic/AutomorphicVarieties/NumberFieldTowers/XTwApplyApply.lean\#L11-L12}{\leanfont Tower.GenLE}}
\expandafter\newcommand\csname LN@GenLEfinite\endcsname{\href{https://github.com/Deicyde/riemann-hypothesis/blob/40e56bab026552e089d7afcfeff08d2defb9c005/src/Development/AutomorphicArithmetic/AutomorphicVarieties/NumberFieldTowers/XTwApplyApply.lean\#L14-L17}{\leanfont Tower.GenLE.finite}}
\expandafter\newcommand\csname LN@cokerLen\endcsname{\href{https://github.com/Deicyde/riemann-hypothesis/blob/40e56bab026552e089d7afcfeff08d2defb9c005/src/Development/AutomorphicArithmetic/AutomorphicVarieties/NumberFieldTowers/XTwApplyApply.lean\#L24-L25}{\leanfont Tower.cokerLen}}
\expandafter\newcommand\csname LN@xRank\endcsname{\href{https://github.com/Deicyde/riemann-hypothesis/blob/40e56bab026552e089d7afcfeff08d2defb9c005/src/Development/AutomorphicArithmetic/AutomorphicVarieties/NumberFieldTowers/XTwApplyApply.lean\#L34-L34}{\leanfont Tower.xRank}}
\expandafter\newcommand\csname LN@IsTower\endcsname{\href{https://github.com/Deicyde/riemann-hypothesis/blob/40e56bab026552e089d7afcfeff08d2defb9c005/src/Development/AutomorphicArithmetic/AutomorphicVarieties/NumberFieldTowers/XTwApplyApply.lean\#L120-L124}{\leanfont Tower.IsTower}}
\expandafter\newcommand\csname LN@IsDesc\endcsname{\href{https://github.com/Deicyde/riemann-hypothesis/blob/40e56bab026552e089d7afcfeff08d2defb9c005/src/Development/AutomorphicArithmetic/AutomorphicVarieties/NumberFieldTowers/IdealNormClassCongr.lean\#L99-L101}{\leanfont Tower.IsDesc}}
\expandafter\newcommand\csname LN@cycMemPrinc\endcsname{\href{https://github.com/Deicyde/riemann-hypothesis/blob/40e56bab026552e089d7afcfeff08d2defb9c005/src/Development/AutomorphicArithmetic/AutomorphicVarieties/NumberFieldTowers/CycMemPrincIff.lean\#L7-L34}{\leanfont cyc\_mem\_princ\_iff}}
\expandafter\newcommand\csname LN@leFinrank\endcsname{\href{https://github.com/Deicyde/riemann-hypothesis/blob/40e56bab026552e089d7afcfeff08d2defb9c005/src/Development/AutomorphicArithmetic/AutomorphicVarieties/NumberFieldTowers/CycMemPrincIff.lean\#L36-L49}{\leanfont le\_finrank}}
\expandafter\newcommand\csname LN@exIff\endcsname{\href{https://github.com/Deicyde/riemann-hypothesis/blob/40e56bab026552e089d7afcfeff08d2defb9c005/src/Development/AutomorphicArithmetic/AutomorphicVarieties/NumberFieldTowers/CycMemPrincIff.lean\#L51-L59}{\leanfont ex\_iff}}
\expandafter\newcommand\csname LN@allIff\endcsname{\href{https://github.com/Deicyde/riemann-hypothesis/blob/40e56bab026552e089d7afcfeff08d2defb9c005/src/Development/AutomorphicArithmetic/AutomorphicVarieties/NumberFieldTowers/CycMemPrincIff.lean\#L61-L79}{\leanfont all\_iff}}
\expandafter\newcommand\csname LN@muQ\endcsname{\href{https://github.com/Deicyde/riemann-hypothesis/blob/40e56bab026552e089d7afcfeff08d2defb9c005/src/Development/AutomorphicArithmetic/AutomorphicVarieties/NumberFieldTowers/CycMemPrincIff.lean\#L81-L104}{\leanfont Tower.muQ}}
\expandafter\newcommand\csname LN@dinfTors\endcsname{\href{https://github.com/Deicyde/riemann-hypothesis/blob/40e56bab026552e089d7afcfeff08d2defb9c005/src/Development/AutomorphicArithmetic/AutomorphicVarieties/NumberFieldTowers/DinfTors.lean\#L7-L86}{\leanfont Tower.dinfTors}}
\expandafter\newcommand\csname LN@existsGen\endcsname{\href{https://github.com/Deicyde/riemann-hypothesis/blob/40e56bab026552e089d7afcfeff08d2defb9c005/src/Development/AutomorphicArithmetic/AutomorphicVarieties/NumberFieldTowers/ExistsGen.lean\#L9-L28}{\leanfont Tower.existsGen}}
\expandafter\newcommand\csname LN@genGal\endcsname{\href{https://github.com/Deicyde/riemann-hypothesis/blob/40e56bab026552e089d7afcfeff08d2defb9c005/src/Development/AutomorphicArithmetic/AutomorphicVarieties/NumberFieldTowers/GenGal.lean\#L9-L33}{\leanfont Tower.genGal}}
\expandafter\newcommand\csname LN@shaDGen\endcsname{\href{https://github.com/Deicyde/riemann-hypothesis/blob/40e56bab026552e089d7afcfeff08d2defb9c005/src/Development/AutomorphicArithmetic/AutomorphicVarieties/NumberFieldTowers/ShaDGen.lean\#L9-L59}{\leanfont Tower.shaDGen}}
\expandafter\newcommand\csname LN@dtVal\endcsname{\href{https://github.com/Deicyde/riemann-hypothesis/blob/40e56bab026552e089d7afcfeff08d2defb9c005/src/Development/AutomorphicArithmetic/AutomorphicVarieties/NumberFieldTowers/DinfTransport.lean\#L8-L12}{\leanfont dt\_val\_lt\_one}}
\expandafter\newcommand\csname LN@dtDec\endcsname{\href{https://github.com/Deicyde/riemann-hypothesis/blob/40e56bab026552e089d7afcfeff08d2defb9c005/src/Development/AutomorphicArithmetic/AutomorphicVarieties/NumberFieldTowers/DinfTransport.lean\#L14-L25}{\leanfont dt\_dec}}
\expandafter\newcommand\csname LN@dtIne\endcsname{\href{https://github.com/Deicyde/riemann-hypothesis/blob/40e56bab026552e089d7afcfeff08d2defb9c005/src/Development/AutomorphicArithmetic/AutomorphicVarieties/NumberFieldTowers/DinfTransport.lean\#L27-L34}{\leanfont dt\_ine}}
\expandafter\newcommand\csname LN@dtMem\endcsname{\href{https://github.com/Deicyde/riemann-hypothesis/blob/40e56bab026552e089d7afcfeff08d2defb9c005/src/Development/AutomorphicArithmetic/AutomorphicVarieties/NumberFieldTowers/DinfTransport.lean\#L36-L48}{\leanfont dt\_mem}}
\expandafter\newcommand\csname LN@dinfTransport\endcsname{\href{https://github.com/Deicyde/riemann-hypothesis/blob/40e56bab026552e089d7afcfeff08d2defb9c005/src/Development/AutomorphicArithmetic/AutomorphicVarieties/NumberFieldTowers/DinfTransport.lean\#L50-L76}{\leanfont Tower.dinfTransport}}
\expandafter\newcommand\csname LN@desc\endcsname{\href{https://github.com/Deicyde/riemann-hypothesis/blob/40e56bab026552e089d7afcfeff08d2defb9c005/src/Development/AutomorphicArithmetic/AutomorphicVarieties/NumberFieldTowers/Imports.lean\#L10-L147}{\leanfont Tower.desc}}
\expandafter\newcommand\csname LN@qzAnn\endcsname{\href{https://github.com/Deicyde/riemann-hypothesis/blob/40e56bab026552e089d7afcfeff08d2defb9c005/src/Development/AutomorphicArithmetic/AutomorphicVarieties/NumberFieldTowers/QzSep.lean\#L8-L45}{\leanfont exists\_qz\_ne\_zero\_annihilators\_of\_ne\_zero}}
\expandafter\newcommand\csname LN@G\endcsname{\href{https://github.com/Deicyde/riemann-hypothesis/blob/40e56bab026552e089d7afcfeff08d2defb9c005/src/Formalism/ConjugateSelfDual/Statement.lean\#L6-L6}{\leanfont G}}
\expandafter\newcommand\csname LN@QZ\endcsname{\href{https://github.com/Deicyde/riemann-hypothesis/blob/40e56bab026552e089d7afcfeff08d2defb9c005/src/Formalism/GaloisCohomology/Statement.lean\#L15-L15}{\leanfont QZ}}
\expandafter\newcommand\csname LN@LData\endcsname{\href{https://github.com/Deicyde/riemann-hypothesis/blob/40e56bab026552e089d7afcfeff08d2defb9c005/src/Formalism/GaloisCohomology/Statement.lean\#L54-L63}{\leanfont LData}}
\expandafter\newcommand\csname LN@qzTw\endcsname{\href{https://github.com/Deicyde/riemann-hypothesis/blob/40e56bab026552e089d7afcfeff08d2defb9c005/src/Formalism/GaloisCohomology/Statement.lean\#L67-L69}{\leanfont qzTw}}
\expandafter\newcommand\csname LN@Sha\endcsname{\href{https://github.com/Deicyde/riemann-hypothesis/blob/40e56bab026552e089d7afcfeff08d2defb9c005/src/Formalism/GaloisCohomology/Statement.lean\#L75-L76}{\leanfont Sha}}
\expandafter\newcommand\csname LN@cycZ\endcsname{\href{https://github.com/Deicyde/riemann-hypothesis/blob/40e56bab026552e089d7afcfeff08d2defb9c005/src/ReciprocityTheory/GaloisValuations/Global.lean\#L90-L92}{\leanfont Gal.cycZ}}
\expandafter\newcommand\csname LN@princ\endcsname{\href{https://github.com/Deicyde/riemann-hypothesis/blob/40e56bab026552e089d7afcfeff08d2defb9c005/src/Development/AutomorphicArithmetic/CyclotomicBounds/ClassNumberBounds/TotallyRealRoots/LeNuIffPowMemPrinc.lean\#L16-L16}{\leanfont princ}}
\expandafter\newcommand\csname LN@levU\endcsname{\href{https://github.com/Deicyde/riemann-hypothesis/blob/40e56bab026552e089d7afcfeff08d2defb9c005/src/Development/AutomorphicArithmetic/AutomorphicVarieties/GaloisLibrary/SelmerTransport/ConjStableShaFam.lean\#L187-L187}{\leanfont Base.levU}}
\expandafter\newcommand\csname LN@cycU\endcsname{\href{https://github.com/Deicyde/riemann-hypothesis/blob/40e56bab026552e089d7afcfeff08d2defb9c005/src/Development/AutomorphicArithmetic/AutomorphicVarieties/GaloisLibrary/SelmerTransport/ConjStableShaFam.lean\#L189-L189}{\leanfont Base.cycU}}
\expandafter\newcommand\csname LN@qExp\endcsname{\href{https://github.com/Deicyde/riemann-hypothesis/blob/40e56bab026552e089d7afcfeff08d2defb9c005/src/Development/AutomorphicArithmetic/AutomorphicVarieties/GaloisLibrary/SelmerTransport/ConjStableShaFam.lean\#L193-L193}{\leanfont Base.qExp}}
\expandafter\newcommand\csname LN@gamU\endcsname{\href{https://github.com/Deicyde/riemann-hypothesis/blob/40e56bab026552e089d7afcfeff08d2defb9c005/src/Development/AutomorphicArithmetic/AutomorphicVarieties/GaloisLibrary/SelmerTransport/ConjStableShaFam.lean\#L195-L195}{\leanfont Base.gamU}}
\expandafter\newcommand\csname LN@galOf\endcsname{\href{https://github.com/Deicyde/riemann-hypothesis/blob/40e56bab026552e089d7afcfeff08d2defb9c005/src/Development/AutomorphicArithmetic/AutomorphicVarieties/GaloisLibrary/SelmerTransport/ConjStableShaFam.lean\#L43-L47}{\leanfont Base.galOf}}
\expandafter\newcommand\csname LN@galEmb\endcsname{\href{https://github.com/Deicyde/riemann-hypothesis/blob/40e56bab026552e089d7afcfeff08d2defb9c005/src/Development/AutomorphicArithmetic/AutomorphicVarieties/GaloisLibrary/SelmerTransport/ConjStableShaFam.lean\#L65-L65}{\leanfont Base.galEmb}}
\expandafter\newcommand\csname LN@galEmbFam\endcsname{\href{https://github.com/Deicyde/riemann-hypothesis/blob/40e56bab026552e089d7afcfeff08d2defb9c005/src/Development/AutomorphicArithmetic/AutomorphicVarieties/GaloisLibrary/SelmerTransport/GalEmbFam.lean\#L5-L30}{\leanfont Base.galEmbFam}}
\expandafter\newcommand\csname LN@galEmbTop\endcsname{\href{https://github.com/Deicyde/riemann-hypothesis/blob/40e56bab026552e089d7afcfeff08d2defb9c005/src/Development/AutomorphicArithmetic/AutomorphicVarieties/GaloisLibrary/SelmerTransport/GalEmbTop.lean\#L32-L69}{\leanfont Base.galEmbTop}}
\expandafter\newcommand\csname LN@Dinf\endcsname{\href{https://github.com/Deicyde/riemann-hypothesis/blob/40e56bab026552e089d7afcfeff08d2defb9c005/src/Development/AutomorphicArithmetic/AutomorphicVarieties/GaloisLibrary/KummerTheory/Mk0IdMapCongr.lean\#L195-L195}{\leanfont Kum.Dinf}}
\expandafter\newcommand\csname LN@Xinf\endcsname{\href{https://github.com/Deicyde/riemann-hypothesis/blob/40e56bab026552e089d7afcfeff08d2defb9c005/src/Development/AutomorphicArithmetic/AutomorphicVarieties/GaloisLibrary/KummerTheory/Mk0IdMapCongr.lean\#L197-L197}{\leanfont Kum.Xinf}}
\expandafter\newcommand\csname LN@funRep\endcsname{\href{https://github.com/Deicyde/riemann-hypothesis/blob/40e56bab026552e089d7afcfeff08d2defb9c005/src/Development/AutomorphicArithmetic/AutomorphicVarieties/GaloisLibrary/KummerTheory/Mk0IdMapCongr.lean\#L121-L124}{\leanfont Kum.funRep}}
\expandafter\newcommand\csname LN@EqHom\endcsname{\href{https://github.com/Deicyde/riemann-hypothesis/blob/40e56bab026552e089d7afcfeff08d2defb9c005/src/Development/AutomorphicArithmetic/AutomorphicVarieties/GaloisLibrary/KummerTheory/Mk0IdMapCongr.lean\#L109-L110}{\leanfont Kum.EqHom}}
\expandafter\newcommand\csname LN@Un\endcsname{\href{https://github.com/Deicyde/riemann-hypothesis/blob/40e56bab026552e089d7afcfeff08d2defb9c005/src/Development/AutomorphicArithmetic/AutomorphicVarieties/GaloisLibrary/KummerTheory/Mk0IdMapCongr.lean\#L186-L186}{\leanfont Kum.Un}}
\expandafter\newcommand\csname LN@Procyc\endcsname{\href{https://github.com/Deicyde/riemann-hypothesis/blob/40e56bab026552e089d7afcfeff08d2defb9c005/src/Development/AutomorphicArithmetic/AutomorphicVarieties/GaloisLibrary/KummerTheory/Mk0IdMapCongr.lean\#L191-L191}{\leanfont Kum.Procyc}}
\expandafter\newcommand\csname LN@resInf\endcsname{\href{https://github.com/Deicyde/riemann-hypothesis/blob/40e56bab026552e089d7afcfeff08d2defb9c005/src/Development/AutomorphicArithmetic/AutomorphicVarieties/GaloisLibrary/KummerTheory/Mk0IdMapCongr.lean\#L214-L217}{\leanfont Kum.resInf}}
\expandafter\newcommand\csname LN@funRepMem\endcsname{\href{https://github.com/Deicyde/riemann-hypothesis/blob/40e56bab026552e089d7afcfeff08d2defb9c005/src/Development/AutomorphicArithmetic/AutomorphicVarieties/GaloisLibrary/KummerTheory/FunRepMem.lean\#L5-L46}{\leanfont Kum.funRepMem}}
\expandafter\newcommand\csname LN@dinfNil\endcsname{\href{https://github.com/Deicyde/riemann-hypothesis/blob/40e56bab026552e089d7afcfeff08d2defb9c005/src/Development/AutomorphicArithmetic/AutomorphicVarieties/GaloisLibrary/KummerTheory/DinfNil.lean\#L62-L235}{\leanfont Kum.dinfNil}}
\expandafter\newcommand\csname LN@kumExistsGen\endcsname{\href{https://github.com/Deicyde/riemann-hypothesis/blob/40e56bab026552e089d7afcfeff08d2defb9c005/src/Development/AutomorphicArithmetic/AutomorphicVarieties/GaloisLibrary/KummerTheory/ExistsGen.lean\#L32-L96}{\leanfont Kum.existsGen}}
\expandafter\newcommand\csname LN@resInfInj\endcsname{\href{https://github.com/Deicyde/riemann-hypothesis/blob/40e56bab026552e089d7afcfeff08d2defb9c005/src/Development/AutomorphicArithmetic/AutomorphicVarieties/GaloisLibrary/KummerTheory/ResInfInj.lean\#L58-L171}{\leanfont Kum.resInfInj}}
\expandafter\newcommand\csname LN@resInfOnto\endcsname{\href{https://github.com/Deicyde/riemann-hypothesis/blob/40e56bab026552e089d7afcfeff08d2defb9c005/src/Development/AutomorphicArithmetic/AutomorphicVarieties/GaloisLibrary/KummerTheory/ResInfOnto.lean\#L102-L166}{\leanfont Kum.resInfOnto}}
\expandafter\newcommand\csname LN@zmodToQZ\endcsname{\href{https://github.com/Deicyde/riemann-hypothesis/blob/40e56bab026552e089d7afcfeff08d2defb9c005/src/Development/AutomorphicArithmetic/AutomorphicVarieties/GaloisLibrary/KummerTheory/ZmodToQZ.lean\#L5-L78}{\leanfont Kum.zmodToQZ}}
\expandafter\newcommand\csname LN@HasMuQ\endcsname{\href{https://github.com/Deicyde/riemann-hypothesis/blob/40e56bab026552e089d7afcfeff08d2defb9c005/src/Development/AutomorphicArithmetic/AutomorphicVarieties/GeneratorCountTuples/LemmaK.lean\#L9-L9}{\leanfont HasMuQ}}
\expandafter\newcommand\csname LN@qzG\endcsname{\href{https://github.com/Deicyde/riemann-hypothesis/blob/40e56bab026552e089d7afcfeff08d2defb9c005/src/Development/AutomorphicArithmetic/AutomorphicVarieties/GaloisLibrary/GaloisCohomologyLength/LevelInertiaLeLevelDec.lean\#L22-L22}{\leanfont qzG}}
\expandafter\newcommand\csname LN@galTransport\endcsname{\href{https://github.com/Deicyde/riemann-hypothesis/blob/40e56bab026552e089d7afcfeff08d2defb9c005/src/Development/AutomorphicArithmetic/AutomorphicVarieties/GaloisLibrary/GaloisCohomologyLength/GalTransport.lean\#L7-L71}{\leanfont galTransport}}
\newcommand{\Pin}{40e56bab}
\newcommand{\PinUrl}{https://github.com/Deicyde/riemann-hypothesis/tree/40e56bab026552e089d7afcfeff08d2defb9c005/src}
\newcommand{\DirFiles}{22}
\newcommand{\DirLines}{1{,}672}
\newcommand{\ConcDecls}{70}
\newcommand{\ConcFiles}{27}
\newcommand{\UnitMembers}{39}
\newcommand{\UnitInDir}{39}
\newcommand{\UnitOutside}{0}
\newcommand{\UnitInNote}{39}
\newcommand{\UnitElsewhere}{0}
\newcommand{\ElsewhereRows}{}
\newcommand{\UnitProofs}{21}
\newcommand{\UnitSteps}{213}
\newcommand{\ConcordanceRows}{\crow{dRep}{\lean{dRep}}{XTwApplyApply.lean}{41--43}
\crow{xRep}{\lean{xRep}}{XTwApplyApply.lean}{45--46}
\crow{xRepApply}{\lean{xRepApply}}{XTwApplyApply.lean}{48--48}
\crow{twist}{\lean{twist}}{XTwApplyApply.lean}{50--54}
\crow{dTw}{\lean{dTw}}{XTwApplyApply.lean}{56--56}
\crow{xTw}{\lean{xTw}}{XTwApplyApply.lean}{58--58}
\crow{dTwApply}{\lean{dTwApply}}{XTwApplyApply.lean}{60--61}
\crow{xTwApply}{\lean{xTwApply}}{XTwApplyApply.lean}{63--64}
\crow{xTwApplyApply}{\lean{xTwApplyApply}}{XTwApplyApply.lean}{66--69}
\crow{ShaD}{\lean{ShaD}}{XTwApplyApply.lean}{73--74}
\crow{IsGen}{\lean{IsGen}}{XTwApplyApply.lean}{36--37}
\crow{cZero}{\lean{cZero}}{XTwApplyApply.lean}{104--104}
\crow{layU}{\lean{layU}}{XTwApplyApply.lean}{108--108}
\crow{layF}{\lean{layF}}{XTwApplyApply.lean}{110--110}
\crow{XN}{\lean{XN}}{XTwApplyApply.lean}{112--112}
\crow{XX}{\lean{XX}}{XTwApplyApply.lean}{114--114}
\crow{XP}{\lean{XP}}{XTwApplyApply.lean}{116--116}
\crow{GenLE}{\lean{GenLE}}{XTwApplyApply.lean}{11--12}
\crow{GenLEfinite}{\lean{GenLEfinite}}{XTwApplyApply.lean}{14--17}
\crow{cokerLen}{\lean{cokerLen}}{XTwApplyApply.lean}{24--25}
\crow{xRank}{\lean{xRank}}{XTwApplyApply.lean}{34--34}
\crow{IsTower}{\lean{IsTower}}{XTwApplyApply.lean}{120--124}
\crow{IsDesc}{\lean{IsDesc}}{IdealNormClassCongr.lean}{99--101}
\crow{cycMemPrinc}{\lean{cycMemPrinc}}{CycMemPrincIff.lean}{7--34}
\crow{leFinrank}{\lean{leFinrank}}{CycMemPrincIff.lean}{36--49}
\crow{exIff}{\lean{exIff}}{CycMemPrincIff.lean}{51--59}
\crow{allIff}{\lean{allIff}}{CycMemPrincIff.lean}{61--79}
\crow{muQ}{\lean{muQ}}{CycMemPrincIff.lean}{81--104}
\crow{dinfTors}{\lean{dinfTors}}{DinfTors.lean}{7--86}
\crow{existsGen}{\lean{existsGen}}{ExistsGen.lean}{9--28}
\crow{genGal}{\lean{genGal}}{GenGal.lean}{9--33}
\crow{shaDGen}{\lean{shaDGen}}{ShaDGen.lean}{9--59}
\crow{dtVal}{\lean{dtVal}}{DinfTransport.lean}{8--12}
\crow{dtDec}{\lean{dtDec}}{DinfTransport.lean}{14--25}
\crow{dtIne}{\lean{dtIne}}{DinfTransport.lean}{27--34}
\crow{dtMem}{\lean{dtMem}}{DinfTransport.lean}{36--48}
\crow{dinfTransport}{\lean{dinfTransport}}{DinfTransport.lean}{50--76}
\crow{desc}{\lean{desc}}{Imports.lean}{10--147}
\crow{qzAnn}{\lean{qzAnn}}{QzSep.lean}{8--45}
\crow{G}{\lean{G}}{Statement.lean}{6--6}
\crow{QZ}{\lean{QZ}}{Statement.lean}{15--15}
\crow{LData}{\lean{LData}}{Statement.lean}{54--63}
\crow{qzTw}{\lean{qzTw}}{Statement.lean}{67--69}
\crow{Sha}{\lean{Sha}}{Statement.lean}{75--76}
\crow{cycZ}{\lean{cycZ}}{Global.lean}{90--92}
\crow{princ}{\lean{princ}}{LeNuIffPowMemPrinc.lean}{16--16}
\crow{levU}{\lean{levU}}{ConjStableShaFam.lean}{187--187}
\crow{cycU}{\lean{cycU}}{ConjStableShaFam.lean}{189--189}
\crow{qExp}{\lean{qExp}}{ConjStableShaFam.lean}{193--193}
\crow{gamU}{\lean{gamU}}{ConjStableShaFam.lean}{195--195}
\crow{galOf}{\lean{galOf}}{ConjStableShaFam.lean}{43--47}
\crow{galEmb}{\lean{galEmb}}{ConjStableShaFam.lean}{65--65}
\crow{galEmbFam}{\lean{galEmbFam}}{GalEmbFam.lean}{5--30}
\crow{galEmbTop}{\lean{galEmbTop}}{GalEmbTop.lean}{32--69}
\crow{Dinf}{\lean{Dinf}}{Mk0IdMapCongr.lean}{195--195}
\crow{Xinf}{\lean{Xinf}}{Mk0IdMapCongr.lean}{197--197}
\crow{funRep}{\lean{funRep}}{Mk0IdMapCongr.lean}{121--124}
\crow{EqHom}{\lean{EqHom}}{Mk0IdMapCongr.lean}{109--110}
\crow{Un}{\lean{Un}}{Mk0IdMapCongr.lean}{186--186}
\crow{Procyc}{\lean{Procyc}}{Mk0IdMapCongr.lean}{191--191}
\crow{resInf}{\lean{resInf}}{Mk0IdMapCongr.lean}{214--217}
\crow{funRepMem}{\lean{funRepMem}}{FunRepMem.lean}{5--46}
\crow{dinfNil}{\lean{dinfNil}}{DinfNil.lean}{62--235}
\crow{kumExistsGen}{\lean{kumExistsGen}}{ExistsGen.lean}{32--96}
\crow{resInfInj}{\lean{resInfInj}}{ResInfInj.lean}{58--171}
\crow{resInfOnto}{\lean{resInfOnto}}{ResInfOnto.lean}{102--166}
\crow{zmodToQZ}{\lean{zmodToQZ}}{ZmodToQZ.lean}{5--78}
\crow{HasMuQ}{\lean{HasMuQ}}{LemmaK.lean}{9--9}
\crow{qzG}{\lean{qzG}}{LevelInertiaLeLevelDec.lean}{22--22}
\crow{galTransport}{\lean{galTransport}}{GalTransport.lean}{7--71}}
\expandafter\newcommand\csname LQ@desc\endcsname{\leanquoteline{0}{theorem UpperTheory.V.Tower.desc :}\leanquoteline{4}{∀ (k L Ω : Type) [Field k] [NumberField k] [Field L] [NumberField L] [Field Ω] [Algebra k L] [Algebra k Ω] [Algebra L Ω] [IsScalarTower k L Ω] [IsAlgClosure k Ω] [IsAlgClosure L Ω] (l : ℕ) [Fact l.Prime] (D : LData L Ω l) (e : ℤ) (U : Subgroup (G k Ω)) [U.Normal], e ≠ 0 → Procyc L Ω l → cycU l U = cycU l (galOf k L Ω) →}\leanquoteline{8}{∃ Φ : Sha D e ≃ₗ[ℤ\_[l]] ShaD l l U (galOf k L Ω) e, IsDesc L D e U Φ.toLinearMap}}
\expandafter\newcommand\csname LQ@shaDGen\endcsname{\leanquoteline{0}{theorem UpperTheory.V.Tower.shaDGen :}\leanquoteline{4}{∀ (k Ω : Type) [Field k] [NumberField k] [Field Ω] [Algebra k Ω] [IsAlgClosure k Ω] (l : ℕ) [Fact l.Prime] (N : ℕ) (U V : Subgroup (G k Ω)) [U.Normal] (e : ℤ) (γ : G k Ω), IsClosed (U : Set (G k Ω)) → V ≤ U → IsGen l V γ →}\leanquoteline{8}{ShaD l N U V e = LinearMap.ker (dTw l N U e γ - 1) ∧ ShaD l N U V e = LinearMap.ker (dTw l N U e γ⁻¹ - 1)}}
\expandafter\newcommand\csname LQ@muQ\endcsname{\leanquoteline{0}{theorem UpperTheory.V.Tower.muQ :}\leanquoteline{4}{∀ (L Ω : Type) [Field L] [NumberField L] [Field Ω] [Algebra L Ω] [IsAlgClosure L Ω] (l : ℕ) [Fact l.Prime], (∀ c : ℕ, (∃ ζ : L, IsPrimitiveRoot ζ (l \textasciicircum{} c)) ↔ c ≤ cZero l L) ∧ (∀ c : ℕ, Un L Ω l c = ⊤ ↔ c ≤ cZero l L) ∧}\leanquoteline{8}{(HasMuQ l L ↔ qExp l ≤ cZero l L) ∧ (HasMuQ l L → IsTotallyComplex L ∧ Procyc L Ω l) ∧ (Procyc L Ω l ↔ l ≠ 2 ∨ HasMuQ l L)}}
% the notations of the functions that the archive defines: one slot for each explicit argument
\newcommand{\dRep}[3]{\rho_{#1,#2,#3}}
\newcommand{\xRep}[3]{\rho^{\vee}_{#1,#2,#3}}
\newcommand{\twist}[2]{#1\otimes#2}
\newcommand{\dTw}[4]{D_\infty(#1,#2,#3)(-#4)}
\newcommand{\dTwOp}[4]{T_{#1,#2,#3;#4}}
\newcommand{\xTw}[4]{X_\infty(#1,#2,#3)(#4)}
\newcommand{\xTwOp}[4]{T^{\vee}_{#1,#2,#3;#4}}
\newcommand{\ShaD}[5]{D_\infty(#1,#2,#3)(-#5)^{#4}}
\newcommand{\cZero}[2]{a_{#1}(#2)}
\newcommand{\layU}[4]{\Gal(#3/#2)_{#1,\,a_{#1}(#2)+#4}}
\newcommand{\layF}[4]{#2\bigl(\mu_{#1^{a_{#1}(#2)+#4}}(#3)\bigr)}
\newcommand{\XN}[4]{X_\infty(#1,#4,\Gal(#3/#2))}
\newcommand{\XX}[2]{X_\infty(#1,1,\Gal(\overline{#2}/#2))}
\newcommand{\XP}[2]{X_\infty(#1,#1,\Gal(\overline{#2}/#2))}
\newcommand{\xRank}[3]{\len D_\infty(#1,#2,#3)[#1]}
\newcommand{\GalG}[2]{\Gal(#2/#1)}
\newcommand{\QZl}[1]{\Q_{#1}/\Z_{#1}}
\newcommand{\LData}[3]{\operatorname{Dat}_{#1}(#3/#2)}
\newcommand{\qzTw}[2]{\mathcal Q(#1)(-#2)}
\newcommand{\ShaOf}[2]{\text{\normalfont Ш}^1(#1,#2)}
\newcommand{\cyc}[4]{\kappa_{#3}^{#2/#1}(#4)}
\newcommand{\princ}[2]{U_{#1}^{(#2)}}
\newcommand{\levU}[3]{#2_{#1,#3}}
\newcommand{\cycU}[2]{#2_{#1,\infty}}
\newcommand{\qExp}[1]{q_{#1}}
\newcommand{\gamU}[2]{#2_{#1,q_{#1}}}
\newcommand{\galOf}[3]{\Gal_{#1}(#3/#2)}
\newcommand{\galEmb}[4]{\iota_{#2/#1}^{#3}(#4)}
\newcommand{\Dinf}[3]{D_\infty(#1,#2,#3)}
\newcommand{\Xinf}[3]{X_\infty(#1,#2,#3)}
\newcommand{\funRep}[3]{\mathcal F_{#1}(#3,#2)}
\newcommand{\EqHom}[3]{\mathcal E_{#1}(#2;#3)}
\newcommand{\Un}[4]{\Gal(#2/#1)_{#3,#4}}
\newcommand{\Procyc}[3]{\mathrm{P}_{#3}(#2/#1)}
\newcommand{\resInf}[3]{\operatorname{res}_{\infty}(#1,#2;#3)}
\newcommand{\qzG}[4]{\mathcal Q_{#3}^{#2/#1}(-#4)}
\newcommand{\Dec}[1]{G_{#1}}
\newcommand{\GS}[2]{G_{#2}(#1)}
\newcommand{\Ine}[1]{I_{#1}}
\newcommand{\XDual}[4]{D_{#1}(#2;#3,#4)}
\newcommand{\vl}[2]{v_{#1}(#2)}
\newcommand{\vlb}[2]{v_{#1}\bigl(#2\bigr)}
\newcommand{\vlB}[2]{v_{#1}\Bigl(#2\Bigr)}


\linespread{1.08}                       % 11 pt on 14.7 pt
\setlength{\parindent}{0pt}
\setlength{\parskip}{5.5pt plus 1.5pt minus 0.5pt}
\setlength{\abovedisplayskip}{7.5pt plus 2pt minus 3pt}
\setlength{\belowdisplayskip}{7.5pt plus 2pt minus 3pt}
\setlength{\abovedisplayshortskip}{3pt plus 2pt}
\setlength{\belowdisplayshortskip}{6pt plus 2pt minus 2pt}
\emergencystretch=1.5em
\clubpenalty=10000 \widowpenalty=10000 \displaywidowpenalty=10000
\relpenalty=7000 \binoppenalty=9000
\raggedbottom
\pagestyle{fancy}\fancyhf{}\renewcommand{\headrulewidth}{0pt}
\fancyfoot[C]{\sffamily\fontsize{7.6}{9}\selectfont\color{inkfaint}\thepage\,/\,\pageref*{LastPage}}
\setlist{topsep=2pt,partopsep=0pt,parsep=0pt,itemsep=4pt,leftmargin=2.1em,labelsep=0.55em}

\newcommand{\Q}{\mathbb{Q}}
\newcommand{\Z}{\mathbb{Z}}
\newcommand{\F}{\mathbb{F}}
\newcommand{\len}{\operatorname{len}}
\newcommand{\cok}{\operatorname{cok}}
\newcommand{\End}{\operatorname{End}}
\newcommand{\Hom}{\operatorname{Hom}}
\newcommand{\Gal}{\operatorname{Gal}}
\newcommand{\diag}{\operatorname{diag}}
\newcommand{\kw}[1]{\emph{#1}}
\newcommand{\m}[1]{\mbox{$#1$}}        % a short formula that must not be broken across lines
\newcommand{\head}[1]{\par\addvspace{17pt plus 4pt minus 2pt}\Needspace{6\baselineskip}%
  \noindent{\fontsize{13.2}{16}\selectfont\bfseries\color{ink}#1}\par\nopagebreak\vspace{-1pt}}
\newcommand{\sect}[2]{\refstepcounter{section}\label{#2}\par\addvspace{19pt plus 4pt minus 2pt}\Needspace{7\baselineskip}%
  \noindent\llap{{\fontsize{13.2}{16}\selectfont\bfseries\color{inkfaint}\thesection}\hspace{9pt}}%
  {\fontsize{13.2}{16}\selectfont\bfseries\color{ink}#1}\par\nopagebreak\vspace{-1pt}}
\newcommand{\subhead}[1]{\par\addvspace{9pt plus 2pt}\Needspace{4\baselineskip}\noindent\textbf{#1}\enspace\ignorespaces}
% where a statement lives in the Lean archive
\newcommand{\leantag}{{\sffamily\bfseries\fontsize{6.3}{8}\selectfont\color{leancol}\addfontfeatures{LetterSpace=7}LEAN}}
\newcommand{\sep}{\unskip\nobreak\hspace{5pt}{\color{inkfaint}\textperiodcentered}\hspace{5pt plus 2pt}\ignorespaces}
\newcommand{\leanline}[1]{\par\nopagebreak\vspace{-2.5pt}{\fontsize{8.9}{12.4}\selectfont\color{inksoft}\raggedright
  \noindent\llap{\leantag\hspace{9pt}}#1\par}\vspace{1pt}}
% a displayed statement in a tinted box: #1 colour, #2 text
\newcommand{\statement}[2]{\par\addvspace{7pt}\begingroup\setlength{\fboxsep}{8.5pt}\setlength{\fboxrule}{0.6pt}%
  \noindent\fcolorbox{#1!30!white}{#1!3!white}{\parbox{\dimexpr\linewidth-2\fboxsep-2\fboxrule\relax}{#2}}%
  \endgroup\par\addvspace{7pt}}
\newcommand{\nt}[2]{\par\noindent\hangindent=1.1em\hangafter=1 #1\enspace{\color{inksoft}#2}\par\vspace{1pt}}
\newcommand{\legdot}[1]{\tikz[baseline=-0.62ex]\fill[#1] (0,0) circle (2.5pt);\hspace{3pt}}
\newcommand{\crow}[4]{\csname CS@#1\endcsname & #2 & \textcolor{inksoft}{#3} & \textcolor{inksoft}{#4}\\}
\newcommand{\leanquoteline}[2]{\par\noindent\hangindent=\dimexpr#1\fontdimen2\font+2.2em\relax\hangafter=1
  \hspace*{#1\fontdimen2\font}#2}
\newcommand{\quoted}[1]{\par\vspace{3pt}{\ttfamily\fontsize{7.5}{10.4}\selectfont\color{leancol}\raggedright
  \setlength{\parskip}{0pt}\leanquote{#1}\par}\vspace{4pt}}
\newcommand{\bx}[3]{{\sffamily\bfseries\fontsize{7.3}{9.2}\selectfont\color{#1!82!black}#2}\\[1.5pt]#3}
\newcommand{\cs}[2]{\expandafter\newcommand\csname CS@#1\endcsname{#2}}

% theorems
\newtheoremstyle{plainx}{7pt plus 2pt}{7pt plus 2pt}{\itshape}{}{\bfseries}{.}{0.5em}{}
\newtheoremstyle{uprightx}{7pt plus 2pt}{7pt plus 2pt}{}{}{\bfseries}{.}{0.5em}{}
\theoremstyle{plainx}
\newtheorem{theorem}{Theorem}[section]
\newtheorem{lemma}[theorem]{Lemma}
\newtheorem{proposition}[theorem]{Proposition}
\newtheorem{corollary}[theorem]{Corollary}
\theoremstyle{uprightx}
\newtheorem{definition}[theorem]{Definition}
\newtheorem{setting}[theorem]{Setting}
\newtheorem{remark}[theorem]{Remark}
\newtheorem{example}[theorem]{Example}
\renewcommand{\qedsymbol}{{\setlength{\fboxsep}{0pt}\setlength{\fboxrule}{0.5pt}\fbox{\rule{0pt}{0.46em}\hspace{0.46em}}}}
\makeatletter
\renewenvironment{proof}[1][Proof]{\par\pushQED{\qed}\normalfont\topsep0pt\partopsep0pt
  \trivlist\item[\hskip\labelsep\itshape #1.]\ignorespaces}{\popQED\endtrivlist\@endpefalse}
\makeatother
% the ambient objects of the note: what the archive passes as implicit and instance arguments
\newenvironment{conventions}{\par\addvspace{5pt}\noindent\textbf{Conventions.}\enspace\ignorespaces}{\par\addvspace{3pt}}

% where each declaration of the concordance is treated
\cs{dRep}{Def.~\ref{d:act}} \cs{xRep}{Def.~\ref{d:act}} \cs{xRepApply}{Def.~\ref{d:act}}
\cs{twist}{Def.~\ref{d:twist}} \cs{dTw}{Def.~\ref{d:twist}} \cs{xTw}{Def.~\ref{d:twist}}
\cs{dTwApply}{Def.~\ref{d:twist}} \cs{xTwApply}{Def.~\ref{d:twist}}
\cs{xTwApplyApply}{Lemma~\ref{l:dual}} \cs{ShaD}{Def.~\ref{d:twist}} \cs{IsGen}{Def.~\ref{d:gen}}
\cs{cZero}{Def.~\ref{d:layers}} \cs{layU}{Def.~\ref{d:layers}} \cs{layF}{Def.~\ref{d:layers}}
\cs{XN}{Def.~\ref{d:char}} \cs{XX}{Def.~\ref{d:char}} \cs{XP}{Def.~\ref{d:char}}
\cs{GenLE}{Def.~\ref{d:tower}} \cs{GenLEfinite}{Def.~\ref{d:tower}}
\cs{cokerLen}{Def.~\ref{d:tower}} \cs{xRank}{Def.~\ref{d:tower}}
\cs{IsTower}{Def.~\ref{d:tower}} \cs{IsDesc}{Def.~\ref{d:desc}}
\cs{cycMemPrinc}{Lemma~\ref{l:fix}} \cs{leFinrank}{Lemma~\ref{l:deg}} \cs{exIff}{Prop.~\ref{p:roots}(a)}
\cs{allIff}{Prop.~\ref{p:roots}(a)} \cs{muQ}{Prop.~\ref{p:roots}} \cs{dinfTors}{Prop.~\ref{p:tors}}
\cs{existsGen}{Prop.~\ref{p:gen}} \cs{genGal}{Prop.~\ref{p:genE}} \cs{shaDGen}{Theorem B}
\cs{dtVal}{Lemma~\ref{l:val}(a)} \cs{dtDec}{Lemma~\ref{l:val}(b)} \cs{dtIne}{Lemma~\ref{l:val}(b)}
\cs{dtMem}{Prop.~\ref{p:transport}, proof} \cs{dinfTransport}{Prop.~\ref{p:transport}} \cs{desc}{Theorem A}
\cs{qzAnn}{Lemma~\ref{l:sep}}
\cs{G}{Def.~\ref{d:G}} \cs{QZ}{Def.~\ref{d:G}} \cs{cycZ}{Def.~\ref{d:G}}
\cs{LData}{Def.~\ref{d:classes}} \cs{qzTw}{Def.~\ref{d:classes}} \cs{Sha}{Def.~\ref{d:classes}}
\cs{princ}{Def.~\ref{d:groups}} \cs{levU}{Def.~\ref{d:groups}} \cs{cycU}{Def.~\ref{d:groups}}
\cs{Un}{Def.~\ref{d:groups}} \cs{qExp}{Def.~\ref{d:groups}} \cs{gamU}{Def.~\ref{d:groups}}
\cs{Procyc}{Def.~\ref{d:groups}} \cs{HasMuQ}{Def.~\ref{d:tower}}
\cs{galOf}{Def.~\ref{d:ext}} \cs{galEmb}{Def.~\ref{d:ext}} \cs{galEmbFam}{Prop.~\ref{p:genE}, proof}
\cs{galEmbTop}{Prop.~\ref{p:genE}, proof} \cs{Dinf}{Def.~\ref{d:char}} \cs{Xinf}{Def.~\ref{d:char}}
\cs{funRep}{Def.~\ref{d:act}} \cs{funRepMem}{Def.~\ref{d:act}} \cs{EqHom}{Def.~\ref{d:eqhom}}
\cs{resInf}{Def.~\ref{d:classes}} \cs{dinfNil}{Theorem B, proof} \cs{kumExistsGen}{Input G}
\cs{resInfInj}{Input R1} \cs{resInfOnto}{Input R2} \cs{zmodToQZ}{Lemma~\ref{l:sep}, proof}
\cs{qzG}{Def.~\ref{d:classes}} \cs{galTransport}{Lemma~\ref{l:theta}}

\begin{document}
\color{ink}

% ------------------------------------------------------------------------------------- title
{\raggedright
{\sffamily\bfseries\fontsize{7.2}{9}\selectfont\color{leancol}\addfontfeatures{LetterSpace=8}DEFORMALIZATION 3\par}
\vspace{5pt}
{\fontsize{21.5}{25.5}\selectfont\bfseries Descent to a layer of the cyclotomic tower\par}
\vspace{5pt}
{\fontsize{12.4}{16}\selectfont\itshape\color{inksoft} One unit of the Lean archive behind the proof of the
Riemann hypothesis announced by L.\,Alpöge, read back into mathematics\par}}
\vspace{8pt}
{\color{rulegrey}\hrule height 0.6pt}
\vspace{7pt}
\thispagestyle{empty}

\textbf{Summary.} Let $F$ be a number field with algebraic closure~$\Omega$, and let $\ell$ be a prime. The group
\m{\GalG{F}{\Omega}} acts by conjugation on the module of the locally constant characters of the kernel of its
$\ell$-adic cyclotomic character with values in~$\QZl{\ell}$, and this action can be twisted by a power of the
cyclotomic character. We prove two statements that the first note of this series~\cite{N1} used as inputs.
First, for a finite extension $E$ of~$F$ inside the cyclotomic tower, restriction identifies the locally trivial
classes of Galois cohomology over~$E$ with the characters fixed by the twisted action of the subgroup that belongs
to~$E$; here two theorems on restriction are taken from a neighbouring unit of the archive. Second, these fixed
elements are the kernel of one operator \m{T-1}, where $T$ is the twisted action of an element whose
cyclotomic value generates the image of the group. On the way we describe which roots of unity of $\ell$-power
order a number field contains in terms of the cyclotomic character.

\statement{leancol}{\fontsize{9.6}{13.2}\selectfont\textbf{What this note is.} On 20 November 2026 L.\,Alpöge
announced a proof of the Riemann hypothesis, with a Lean archive \cite{Al}. That is a claim, and nothing here asserts
that the hypothesis is proved. The reading project divides the part of the archive that the final theorems use
into units. This note treats one unit, on the layers of the cyclotomic tower and the twisted actions: all
\UnitMembers{} of its declarations that the final theorems use. Three results are taken from other units and
stated as Inputs G, R1 and R2; they are not proved here. Everything else is proved in full. Every statement was
compared with its formal counterpart at commit \href{\PinUrl}{{\leanfont\Pin}} of the public mirror; Appendix~A is
the concordance of \ConcDecls{} declarations. The proofs were written from the structured proofs of the reading
project (§\ref{s:notes}).}

% ------------------------------------------------------------------------------------- 1
\sect{The statements}{s:intro}
The first note~\cite{N1} bounded the lengths of groups of locally trivial Galois cohomology along the cyclotomic
tower of a number field. It reduced them to cokernels of one operator on an Iwasawa module, and for this it needed
two facts about the dual module, its Inputs 2a and~2b. Both are theorems of the unit of the archive treated here.
The notation is set up in §\ref{s:setting}, and every function that the archive defines is written with all the
arguments of its formal definition. In particular, for a prime~$\ell$, a number field~$F$ with algebraic
closure~$\Omega$, a subgroup $H$ of~$\GalG{F}{\Omega}$ and integers \m{N\ge0} and~$h$, the group $\cycU{\ell}{H}$ is
the intersection of $H$ with the kernel of the cyclotomic character~$\cyc{F}{\Omega}{\ell}{\cdot}$; the module
$\Dinf{\ell}{N}{H}$ consists of characters of~$\cycU{\ell}{H}$; the twisted action of \m{g\in\GalG{F}{\Omega}} on
it, when $H$ is normal, is
\[
\dTwOp{\ell}{N}{H}{h}(g)\,\varphi=\cyc{F}{\Omega}{\ell}{g}^{-h}\,\dRep{\ell}{N}{H}(g)\,\varphi ,
\]
where $\dRep{\ell}{N}{H}(g)$ is conjugation by~$g$; and $\ShaD{\ell}{N}{H}{V}{h}$ is the module of the elements
fixed by $\dTwOp{\ell}{N}{H}{h}(g)$ for all $g$ in a subgroup~$V$.

\statement{arith}{\phantomsection\label{t:A}\textbf{Theorem A} (descent).\enspace\emph{In Setting~\ref{set:ext}, let
\m{\mathcal D\in\LData{\ell}{E}{\Omega}} be a datum for the field~$E$ (Definition~\ref{d:classes}), let \m{h\ne0} be
an integer, and assume that condition~$\Procyc{E}{\Omega}{\ell}$ holds: $\ell$ is odd, or
$\cyc{E}{\Omega}{\ell}{\GalG{E}{\Omega}}\subseteq\princ{\ell}{2}$. Let $H$ be a normal subgroup of~$\GalG{F}{\Omega}$
with \m{\cycU{\ell}{H}=\cycU{\ell}{\galOf{F}{E}{\Omega}}}. Then there is an isomorphism of\/ $\Z_\ell$-modules}
\[
\Phi\colon\ \ShaOf{\mathcal D}{h}\ \longrightarrow\ \ShaD{\ell}{\ell}{H}{\galOf{F}{E}{\Omega}}{h}
\]
\emph{that is given by restriction for $E$, $\mathcal D$, $h$ and~$H$ (Definition~\ref{d:desc}):
\m{(\Phi\xi)(u')=\resInf{\mathcal D}{h}{\xi}(u)} for all \m{\xi\in\ShaOf{\mathcal D}{h}}, all
\m{u\in\cycU{\ell}{\GalG{E}{\Omega}}} and all \m{u'\in\cycU{\ell}{H}} with \m{u'=\galEmb{F}{E}{\Omega}{u}}.}

\vspace{3pt}\phantomsection\label{t:B}\textbf{Theorem B} (invariants).\enspace\emph{In Setting~\ref{set:base}, let
$H$ be a closed normal subgroup of~$\GalG{F}{\Omega}$, let \m{N\ge0} and $h$ be integers, and let $V$ be a
subgroup of~$H$. Let $\gamma$ be an $\ell$-generator for~$V$ (Definition~\ref{d:gen}): \m{\gamma\in V}, and
$\cyc{F}{\Omega}{\ell}{V}$ lies in the closure of the set of powers \m{\cyc{F}{\Omega}{\ell}{\gamma}^j}, \m{j\in\Z}. Then}
\[
\ShaD{\ell}{N}{H}{V}{h}=\ker\bigl(\dTwOp{\ell}{N}{H}{h}(\gamma)-1\bigr)=\ker\bigl(\dTwOp{\ell}{N}{H}{h}(\gamma^{-1})-1\bigr) .
\]}
\leanline{A \lean{desc}, in the form \lean{IsDesc}\sep B \lean{shaDGen}}

In Theorem~A the fields, the prime, the datum~$\mathcal D$, the integer~$h$ and the group~$H$ are given first, and
the isomorphism~$\Phi$ is asserted for them. By Proposition~\ref{p:genE}(a) the map
$\galEmb{F}{E}{\Omega}{\cdot}$, which regards an automorphism of $\Omega$ over~$E$ as one over~$F$, carries
$\cycU{\ell}{\GalG{E}{\Omega}}$ bijectively onto \m{\cycU{\ell}{\galOf{F}{E}{\Omega}}=\cycU{\ell}{H}}. So $\Phi\xi$ is
the map $\resInf{\mathcal D}{h}{\xi}$ read on~$\cycU{\ell}{H}$: the theorem says that restriction takes the locally
trivial classes into the invariants, and that it is an isomorphism onto them.

The data~$\mathcal D$, the group $\ShaOf{\mathcal D}{h}$ of locally trivial classes and the restriction map are
recalled in §\ref{s:proofA}. In Theorem~A the case \m{H=\GalG{F}{\Omega}} is the one used in~\cite{N1}: there the
hypothesis on $H$ says that \m{\cycU{\ell}{\GalG{F}{\Omega}}\subseteq\galOf{F}{E}{\Omega}}, that is, that $E$ lies in
the field generated over $F$ by all roots of unity of $\ell$-power order in~$\Omega$. Theorem~A rests on two theorems
on the restriction map, Inputs R1 and~R2 of §\ref{s:proofA}; they carry its arithmetic content and belong to
another unit. Theorem~B is proved in full, from the local constancy established in the second note~\cite{N2}.

% ------------------------------------------------------------------------------------- 2
\sect{The setting}{s:setting}
\begin{conventions}\phantomsection\label{c:conventions}
Rings are commutative, and a module is a module over the ring that is named. A \kw{representation} of a
group~$\Gamma$ on a module $M$ over a ring~$\Lambda$ is a homomorphism from $\Gamma$ to the monoid of the
$\Lambda$-linear maps \m{M\to M}. For a prime number~$\ell$ we write $\Z_\ell$ for the ring of $\ell$-adic integers,
$\Q_\ell$ for its field of fractions and $\Z_\ell^\times$ for the group of units of~$\Z_\ell$ with its $\ell$-adic
topology. A \kw{number field} is a field of finite degree over~$\Q$. When a field~$\Omega$ is called an algebraic
closure of a field~$F$, an embedding of $F$ into~$\Omega$ belongs to it, and $F$ is regarded as a subfield
of~$\Omega$. A \kw{finite extension of $F$ in~$\Omega$} is a subfield $E$ of~$\Omega$ that contains~$F$ and has finite
degree over it; the embeddings of $F$ into~$E$ and of $E$ into~$\Omega$ are compatible with the one of $F$
into~$\Omega$. A subgroup of a topological group carries the induced topology. The \kw{length} $\len M$ of a
module~$M$ is the supremum of the lengths of its chains of submodules, a non-negative integer or~$\infty$.
\end{conventions}

A statement of this note either states all its hypotheses itself, or it cites one of three Settings that hold
them: the next one, Setting~\ref{set:ext}, and Setting~\ref{set:tr} of §\ref{s:transport}.

\begin{setting}\label{set:base}
$\ell$ is a prime number, $F$ is a number field, and $\Omega$ is an algebraic closure of~$F$.
\end{setting}

\begin{definition}[the group, the coefficients, the character]\label{d:G}
In Setting~\ref{set:base}, $\GalG{F}{\Omega}$ is the group of the automorphisms of the field~$\Omega$ that fix every
element of~$F$, with its usual topology, in which the subgroups of the elements that fix a finite extension of
$F$ in~$\Omega$ are a basis of neighbourhoods of~$1$. This group is compact \cite[Ch.\,I, §1]{NSW}. The group
$\QZl{\ell}$ is the quotient of $\Q_\ell$ by~$\Z_\ell$, a module over~$\Z_\ell$. The \kw{$\ell$-adic cyclotomic
character} of $\Omega$ over~$F$ is the continuous homomorphism
\[
\cyc{F}{\Omega}{\ell}{\cdot}\colon\ \GalG{F}{\Omega}\longrightarrow\Z_\ell^\times\qquad\text{with}\qquad
g(\zeta)=\zeta^{\cyc{F}{\Omega}{\ell}{g}}
\]
for every \m{g\in\GalG{F}{\Omega}} and every root of unity $\zeta$ of $\ell$-power order in~$\Omega$. For a subgroup
$H$ of~$\GalG{F}{\Omega}$ we write $\cyc{F}{\Omega}{\ell}{H}$ for its image, the group of the values
$\cyc{F}{\Omega}{\ell}{g}$ with \m{g\in H}.
\end{definition}
\leanline{\lean{G}\sep\lean{QZ}\sep\lean{cycZ}}

\begin{definition}[levels]\label{d:groups}
In Setting~\ref{set:base}, let \m{n\ge0} be an integer and $H$ a subgroup of~$\GalG{F}{\Omega}$.
\begin{enumerate}[label=\textup{(\roman*)},topsep=3pt,itemsep=1pt]
\item $\princ{\ell}{n}$ is the group of the units of $\Z_\ell$ that are congruent to~$1$ modulo~$\ell^n$. So
\m{\princ{\ell}{n}=1+\ell^n\Z_\ell} for \m{n\ge1}, and \m{\princ{\ell}{0}=\Z_\ell^\times}.
\item The subgroups of $H$ of level~$n$ and of infinite level are
\[
\levU{\ell}{H}{n}=\{g\in H:\ \cyc{F}{\Omega}{\ell}{g}\in\princ{\ell}{n}\},\qquad
\cycU{\ell}{H}=\{g\in H:\ \cyc{F}{\Omega}{\ell}{g}=1\} .
\]
For the whole group, $\Un{F}{\Omega}{\ell}{n}$ consists of all \m{g\in\GalG{F}{\Omega}} with
\m{\cyc{F}{\Omega}{\ell}{g}\in\princ{\ell}{n}}, and $\cycU{\ell}{\GalG{F}{\Omega}}$ is the kernel
of~$\cyc{F}{\Omega}{\ell}{\cdot}$.
\item \m{\qExp{\ell}=2} if \m{\ell=2}, and \m{\qExp{\ell}=1} if $\ell$ is odd, so that $\ell^{\qExp{\ell}}$ is $4$
or~$\ell$. The group $\gamU{\ell}{H}$ is the subgroup of $H$ of level~$\qExp{\ell}$.
\item \kw{Condition~$\Procyc{F}{\Omega}{\ell}$ holds} if $\ell$ is odd or
\m{\cyc{F}{\Omega}{\ell}{\GalG{F}{\Omega}}\subseteq\princ{\ell}{2}}. For \m{\ell=2} this says that
$\Un{F}{\Omega}{\ell}{\qExp{\ell}}$ is the whole group.
\end{enumerate}
\end{definition}
\leanline{\lean{princ}\sep\lean{levU}\sep\lean{cycU}\sep\lean{Un}\sep\lean{qExp}\sep\lean{gamU}\sep\lean{Procyc}}

\begin{setting}\label{set:ext}
$\ell$ is a prime number, $F$ is a number field, $\Omega$ is an algebraic closure of~$F$, and $E$ is a finite
extension of $F$ in~$\Omega$. Then $E$ is a number field and $\Omega$ is an algebraic closure of~$E$, so that
Setting~\ref{set:base} holds for~$E$ in the place of~$F$ as well.
\end{setting}

\begin{definition}[the subgroup of an extension]\label{d:ext}
In Setting~\ref{set:ext}, $\galOf{F}{E}{\Omega}$ is the subgroup of the \m{g\in\GalG{F}{\Omega}} that fix every
element of~$E$. An element $g$ of~$\GalG{E}{\Omega}$ fixes every element of~$F$; regarded as an element
of~$\GalG{F}{\Omega}$ it is written~$\galEmb{F}{E}{\Omega}{g}$. The map
\m{\galEmb{F}{E}{\Omega}{\cdot}\colon\GalG{E}{\Omega}\to\GalG{F}{\Omega}} is an injective homomorphism, and its image
is~$\galOf{F}{E}{\Omega}$.
\end{definition}
\leanline{\lean{galOf}\sep\lean{galEmb}}

\begin{definition}[characters]\label{d:char}
In Setting~\ref{set:base}, a \kw{valuation ring} of $\Omega$ is a subring $A$ such that every \m{t\ne0} in $\Omega$
has $t$ or $t^{-1}$ in~$A$. Its non-units form its maximal ideal~$\mathfrak m_A$. The \kw{decomposition group} of
$A$ is \m{\Dec{A}=\{g\in\GalG{F}{\Omega}:\ g(A)=A\}}, and the \kw{inertia group} $\Ine{A}$ consists of the
\m{g\in\Dec{A}} with \m{g(t)-t\in\mathfrak m_A} for all \m{t\in A}. Both are subgroups of~$\GalG{F}{\Omega}$ for the
fields of the Setting: the notation shows the valuation ring only, and where two fields are in play the text says
in which group they are formed.

Let $C$ be a subgroup of~$\GalG{F}{\Omega}$, let \m{N\ge0} be an integer, $\Lambda$ a ring and $M$ a
$\Lambda$-module. As in the definition of the modules of characters in the second note \cite[§2]{N2},
$\XDual{N}{C}{\Lambda}{M}$ denotes the $\Lambda$-module of the maps
\m{\varphi\colon C\to M} that are locally constant, are additive, \m{\varphi(uw)=\varphi(u)+\varphi(w)}, vanish on
\m{C\cap\Ine{A}} for every valuation ring~$A$ of~$\Omega$, and vanish on \m{C\cap\Dec{A}} for every
valuation ring $A$ with \m{N\in\mathfrak m_A}. For a subgroup $H$ of~$\GalG{F}{\Omega}$ we put
\[
\Dinf{\ell}{N}{H}=\XDual{N}{\cycU{\ell}{H}}{\Z_\ell}{\QZl{\ell}},\qquad
\Xinf{\ell}{N}{H}=\Hom_{\Z_\ell}\bigl(\Dinf{\ell}{N}{H},\QZl{\ell}\bigr) .
\]
The prime, the integer~$N$ and the group~$H$ are written every time. When $H$ is the whole group, the second
module is~$\XN{\ell}{F}{\Omega}{N}$. Let $\overline F$ be the algebraic closure of~$F$ that the library of the
archive constructs; then $\XX{\ell}{F}$ and $\XP{\ell}{F}$ are the modules with \m{N=1} and with \m{N=\ell} for
the whole group over~$\overline F$. The modules of \cite{N1} are those of the whole group with \m{N=\ell}.
\end{definition}
\leanline{\lean{Dinf}\sep\lean{Xinf}\sep for the whole group: \lean{XN}, and its cases \m{N=1} and \m{N=\ell} over
the algebraic closure that the library provides, \lean{XX} and \lean{XP}}

\begin{definition}[actions]\label{d:act}
In Setting~\ref{set:base}, let $C$ be a normal subgroup of~$\GalG{F}{\Omega}$, $\Lambda$ a ring and $M$ a
$\Lambda$-module. The group $\GalG{F}{\Omega}$ acts by conjugation on the $\Lambda$-module of all maps
\m{\varphi\colon C\to M}: the representation $\funRep{\Lambda}{M}{C}$ is given by
\[
\bigl(\funRep{\Lambda}{M}{C}(g)\,\varphi\bigr)(u)=\varphi(g^{-1}ug)\qquad\text{for }g\in\GalG{F}{\Omega},\ u\in C .
\]
Let $H$ be a normal subgroup of~$\GalG{F}{\Omega}$ and \m{N\ge0} an integer. Then $\cycU{\ell}{H}$ is normal
in~$\GalG{F}{\Omega}$, and by the proposition on conjugation of the second note \cite[§2]{N2} the operators
$\funRep{\Z_\ell}{\QZl{\ell}}{\cycU{\ell}{H}}(g)$ map $\Dinf{\ell}{N}{H}$ into itself. Their restrictions to this
module form the representation $\dRep{\ell}{N}{H}$ of $\GalG{F}{\Omega}$ on~$\Dinf{\ell}{N}{H}$; it is
$\Z_\ell$-linear, and
\[
\bigl(\dRep{\ell}{N}{H}(g)\,\varphi\bigr)(u)=\varphi(g^{-1}ug)\qquad\text{for }\varphi\in\Dinf{\ell}{N}{H},\ u\in\cycU{\ell}{H} .
\]
On $\Xinf{\ell}{N}{H}$ the group acts by the representation $\xRep{\ell}{N}{H}$ with
\[
\bigl(\xRep{\ell}{N}{H}(g)\,x\bigr)(\varphi)=x\bigl(\dRep{\ell}{N}{H}(g^{-1})\,\varphi\bigr)\qquad
\text{for }x\in\Xinf{\ell}{N}{H},\ \varphi\in\Dinf{\ell}{N}{H} .
\]
\end{definition}
\leanline{\lean{funRep}\sep\lean{dRep}, a restriction of \lean{funRep} by \lean{funRepMem}\sep\lean{xRep},
\lean{xRepApply}}

\begin{definition}[twists and invariants]\label{d:twist}
Let $\Lambda$ be a ring, $\Gamma$ a group, $M$ a $\Lambda$-module, $\pi$ a representation of $\Gamma$ on~$M$, and
$\omega$ a homomorphism from $\Gamma$ to the group of units of~$\Lambda$. Scalars commute with $\Lambda$-linear
maps, so \m{g\mapsto\omega(g)\,\pi(g)} is again a representation of $\Gamma$ on~$M$. It is the \kw{twist}
$\twist{\pi}{\omega}$ of $\pi$ by~$\omega$.

In Setting~\ref{set:base}, let $H$ be a normal subgroup of~$\GalG{F}{\Omega}$, and let \m{N\ge0} and $h$ be
integers. The cyclotomic character $\cyc{F}{\Omega}{\ell}{\cdot}$ is a homomorphism to the units of~$\Z_\ell$, and
so are its powers. We put
\[
\dTwOp{\ell}{N}{H}{h}=\twist{\dRep{\ell}{N}{H}}{\cyc{F}{\Omega}{\ell}{\cdot}^{-h}},\qquad
\xTwOp{\ell}{N}{H}{h}=\twist{\xRep{\ell}{N}{H}}{\cyc{F}{\Omega}{\ell}{\cdot}^{h}} ,
\]
so that for \m{g\in\GalG{F}{\Omega}}
\[
\dTwOp{\ell}{N}{H}{h}(g)=\cyc{F}{\Omega}{\ell}{g}^{-h}\,\dRep{\ell}{N}{H}(g),\qquad
\xTwOp{\ell}{N}{H}{h}(g)=\cyc{F}{\Omega}{\ell}{g}^{h}\,\xRep{\ell}{N}{H}(g) .
\]
We write $\dTw{\ell}{N}{H}{h}$ for the module $\Dinf{\ell}{N}{H}$ with the action~$\dTwOp{\ell}{N}{H}{h}$ of
$\GalG{F}{\Omega}$, and $\xTw{\ell}{N}{H}{h}$ for the module $\Xinf{\ell}{N}{H}$ with the
action~$\xTwOp{\ell}{N}{H}{h}$. For a subgroup $V$ of~$\GalG{F}{\Omega}$ the module of invariants is
\[
\ShaD{\ell}{N}{H}{V}{h}=\bigl\{\varphi\in\Dinf{\ell}{N}{H}:\ \dTwOp{\ell}{N}{H}{h}(g)\,\varphi=\varphi\ \text{for all }g\in V\bigr\}.
\]
\end{definition}
\leanline{\lean{twist}\sep\lean{dTw}, \lean{dTwApply}\sep\lean{xTw}, \lean{xTwApply}\sep\lean{ShaD}}

\begin{lemma}\label{l:dual}
In Setting~\ref{set:base}, let $H$ be a normal subgroup of~$\GalG{F}{\Omega}$, let \m{N\ge0} and $h$ be integers,
and let \m{g\in\GalG{F}{\Omega}}, \m{x\in\Xinf{\ell}{N}{H}} and \m{\varphi\in\Dinf{\ell}{N}{H}}. Then
\[
\bigl(\xTwOp{\ell}{N}{H}{h}(g)\,x\bigr)(\varphi)=x\bigl(\dTwOp{\ell}{N}{H}{h}(g^{-1})\,\varphi\bigr) .
\]
\end{lemma}
\begin{proof}
The left side is \m{\cyc{F}{\Omega}{\ell}{g}^{h}\,x\bigl(\dRep{\ell}{N}{H}(g^{-1})\,\varphi\bigr)} by
Definitions \ref{d:act} and~\ref{d:twist}. Since
\m{\cyc{F}{\Omega}{\ell}{g}^h=\cyc{F}{\Omega}{\ell}{g^{-1}}^{-h}} and $x$ is $\Z_\ell$-linear, this is the right side.
\end{proof}
\leanline{\lean{xTwApplyApply}}

\begin{definition}\label{d:gen}
In Setting~\ref{set:base}, let $V$ be a subgroup of~$\GalG{F}{\Omega}$. An element \m{\gamma\in V} is an
\kw{$\ell$-generator for~$V$} if $\cyc{F}{\Omega}{\ell}{V}$ is contained in the closure of the set
\m{\{\cyc{F}{\Omega}{\ell}{\gamma}^j:\ j\in\Z\}} in~$\Z_\ell^\times$.
\end{definition}
\leanline{\lean{IsGen}}

\begin{definition}[layers]\label{d:layers}
In Setting~\ref{set:base}, $\cZero{\ell}{F}$ is the largest integer \m{c\ge0} such that $F$ contains a primitive
root of unity of order~$\ell^c$; it exists by Lemma~\ref{l:deg}. For an integer \m{n\ge0} let
$\mu_{\ell^n}(\Omega)$ be the set of the \m{t\in\Omega} with \m{t^{\ell^n}=1}. For \m{m\ge0} the \kw{layer} of
index~$m$ of $F$ in~$\Omega$ is the field
\[
\layF{\ell}{F}{\Omega}{m}
\]
generated over $F$ by the set $\mu_{\ell^n}(\Omega)$ for \m{n=\cZero{\ell}{F}+m}, that is, by the roots of unity
in~$\Omega$ whose order divides this power of~$\ell$. Its group is $\layU{\ell}{F}{\Omega}{m}$ (Lemma~\ref{l:fix}).
\end{definition}
\leanline{\lean{cZero}\sep\lean{layU}\sep\lean{layF}}

\begin{definition}[the tower condition]\label{d:tower}
The main theorem of~\cite{N1} is stated with the following notions.
\begin{enumerate}[label=\textup{(\roman*)},topsep=3pt,itemsep=1pt]
\item Let $\Lambda$ be a ring, $M$ a $\Lambda$-module and \m{n\ge0} an integer. The module $M$ is \kw{generated
by at most $n$ elements over~$\Lambda$} if it is the image of a $\Lambda$-linear map from the free module of
rank~$n$; it is then finitely generated.
\item The archive measures a $\Lambda$-linear map \m{\alpha\colon M\to M} by the length of its cokernel
$M/\alpha(M)$ as a $\Lambda$-module, a non-negative integer or~$\infty$.
\item In Setting~\ref{set:base}, let $H$ be a subgroup of~$\GalG{F}{\Omega}$ and \m{N\ge0} an integer. The
elements of $\Dinf{\ell}{N}{H}$ that are killed by~$\ell$ form the submodule $\Dinf{\ell}{N}{H}[\ell]$, and its
length as a $\Z_\ell$-module is written~$\xRank{\ell}{N}{H}$.
\item Let $\ell$ be a prime number, $F$ a number field and $R$ an integer. The \kw{tower condition holds for
$\ell$, $F$ and~$R$} if the field $F$ is totally complex, it contains a primitive root of unity of
order~$\ell^{\qExp{\ell}}$, \m{R\ge1}, and the module $\XP{\ell}{F}$ is generated by at most $R$ elements
over~$\Z_\ell$. This is the standing hypothesis of that theorem.
\end{enumerate}
\end{definition}
\leanline{\lean{GenLE}, \lean{GenLEfinite}\sep\lean{cokerLen}\sep\lean{xRank}\sep\lean{IsTower}, with
\lean{HasMuQ} for the root of unity of order~$\ell^{\qExp{\ell}}$}

% ------------------------------------------------------------------------------------- 3
\sect{Roots of unity and levels}{s:roots}
\begin{lemma}\label{l:fix}
In Setting~\ref{set:base}, let \m{g\in\GalG{F}{\Omega}} and let \m{n\ge0} be an integer. Then
\m{\cyc{F}{\Omega}{\ell}{g}\in\princ{\ell}{n}} if and only if \m{g(t)=t} for every \m{t\in\Omega} with
\m{t^{\ell^n}=1}.
\end{lemma}
\begin{proof}
For \m{n=0} both conditions hold. Let \m{n\ge1}. The field $\Omega$ is algebraically closed of characteristic
zero, so the \m{t\in\Omega} with \m{t^{\ell^n}=1} form a cyclic group of order~$\ell^n$; let $\zeta$ be a generator.
By the definition of the cyclotomic character we have \m{g(t)=t^{\cyc{F}{\Omega}{\ell}{g}}} for these~$t$, where
the exponent matters only modulo~$\ell^n$. If \m{\cyc{F}{\Omega}{\ell}{g}\equiv1} modulo~$\ell^n$, then \m{g(t)=t}.
Conversely, \m{g(\zeta)=\zeta} gives \m{\zeta^{\cyc{F}{\Omega}{\ell}{g}-1}=1}, so $\ell^n$ divides
\m{\cyc{F}{\Omega}{\ell}{g}-1}.
\end{proof}
\leanline{\lean{cycMemPrinc}}

\begin{lemma}\label{l:deg}
Let $\ell$ be a prime number, $F$ a number field and \m{c\ge0} an integer. If $F$ contains a primitive root of unity
of order~$\ell^c$, then \m{c\le[F:\Q]}.
\end{lemma}
\begin{proof}
Let \m{c\ge1}. The minimal polynomial over $\Q$ of such a root is the cyclotomic polynomial of order~$\ell^c$,
of degree \m{\ell^{c-1}(\ell-1)} \cite[Thm.\,2.5]{Wa}. This degree is at most~$[F:\Q]$, and
\m{c\le\ell^{c-1}\le\ell^{c-1}(\ell-1)}.
\end{proof}
\leanline{\lean{leFinrank}}

\begin{proposition}\label{p:roots}
In Setting~\ref{set:base} the following hold.
\begin{enumerate}[label=\textup{(\alph*)},topsep=3pt,itemsep=1pt]
\item For every integer \m{c\ge0} the following are equivalent: $F$ contains a primitive root of unity of
order~$\ell^c$; \m{c\le\cZero{\ell}{F}}; \m{\cyc{F}{\Omega}{\ell}{\GalG{F}{\Omega}}\subseteq\princ{\ell}{c}}, that
is, $\Un{F}{\Omega}{\ell}{c}$ is the whole group~$\GalG{F}{\Omega}$.
\item $F$ contains a primitive root of unity of order~$\ell^{\qExp{\ell}}$ if and only if
\m{\qExp{\ell}\le\cZero{\ell}{F}}. In that case $F$ is totally complex and
condition~$\Procyc{F}{\Omega}{\ell}$ holds.
\item Condition~$\Procyc{F}{\Omega}{\ell}$ holds if and only if $\ell$ is odd or $F$ contains a primitive root of
unity of order~$\ell^{\qExp{\ell}}$, which for \m{\ell=2} is a primitive fourth root of unity.
\end{enumerate}
\end{proposition}
\begin{proof}
(a) Write \m{a=\cZero{\ell}{F}}. The integers $c$ for which $F$ contains a primitive root of unity of order
$\ell^c$ include~$0$ and are bounded by Lemma~\ref{l:deg}, so there is a largest one, which is~$a$. If $\zeta$ is
such a root of order $\ell^a$ and \m{c\le a}, then $\zeta^{\ell^{a-c}}$ has order~$\ell^c$. This gives the first
equivalence. For the second, suppose \m{\cyc{F}{\Omega}{\ell}{\GalG{F}{\Omega}}\subseteq\princ{\ell}{c}}, and let
$\zeta$ be a primitive root of unity of order $\ell^c$ in~$\Omega$. By Lemma~\ref{l:fix} it is fixed
by~$\GalG{F}{\Omega}$, and the field of the elements fixed by $\GalG{F}{\Omega}$ is~$F$, so \m{\zeta\in F}.
Conversely, if $F$ contains a primitive root $\zeta$ of order~$\ell^c$, every \m{t\in\Omega} with \m{t^{\ell^c}=1}
is a power of~$\zeta$, hence fixed by~$\GalG{F}{\Omega}$, and Lemma~\ref{l:fix} gives
\m{\cyc{F}{\Omega}{\ell}{\GalG{F}{\Omega}}\subseteq\princ{\ell}{c}}.

(b) The first assertion is (a) for \m{c=\qExp{\ell}}. A root of unity of order \m{\ell^{\qExp{\ell}}>2} is not real,
so $F$ has no embedding into the real numbers. If \m{\ell=2}, then \m{\qExp{\ell}=2} and (a) gives
\m{\cyc{F}{\Omega}{\ell}{\GalG{F}{\Omega}}\subseteq\princ{\ell}{2}}.

(c) For odd $\ell$ both statements hold. For \m{\ell=2} condition~$\Procyc{F}{\Omega}{\ell}$ says
\m{\cyc{F}{\Omega}{\ell}{\GalG{F}{\Omega}}\subseteq\princ{\ell}{2}}, which by~(a) means that $F$ contains a
primitive fourth root of unity.
\end{proof}
\leanline{\lean{muQ}\sep (a) \lean{exIff}, \lean{allIff}\sep (b) the root of order~$\ell^{\qExp{\ell}}$:
\lean{HasMuQ}}

% ------------------------------------------------------------------------------------- 4
\sect{Torsion, and elements that act trivially}{s:torsion}
\begin{proposition}\label{p:tors}
In Setting~\ref{set:base}, let $H$ be a subgroup of~$\GalG{F}{\Omega}$ and \m{N\ge0} an integer.
\begin{enumerate}[label=\textup{(\alph*)},topsep=3pt,itemsep=1pt]
\item If $H$ is closed in~$\GalG{F}{\Omega}$, every element of $\Dinf{\ell}{N}{H}$ is killed by a power of~$\ell$.
\item \m{\Dinf{\ell}{N}{H}\subseteq\Dinf{\ell}{1}{H}}, and \m{\xRank{\ell}{N}{H}\le\xRank{\ell}{1}{H}}: the
submodule of the elements killed by~$\ell$ has at most the length of the corresponding submodule
of~$\Dinf{\ell}{1}{H}$.
\item If $H$ is normal in~$\GalG{F}{\Omega}$ and \m{g\in\cycU{\ell}{H}}, then $\dRep{\ell}{N}{H}(g)$ is the identity
of~$\Dinf{\ell}{N}{H}$ and $\xRep{\ell}{N}{H}(g)$ is the identity of~$\Xinf{\ell}{N}{H}$, and so are
$\dTwOp{\ell}{N}{H}{h}(g)$ and $\xTwOp{\ell}{N}{H}{h}(g)$ for every integer~$h$.
\end{enumerate}
\end{proposition}
\begin{proof}
(a) An element of $\QZl{\ell}$ is the class of some \m{s\in\Q_\ell}, and \m{\ell^rs\in\Z_\ell} for large~$r$; so
every element of $\QZl{\ell}$ is killed by a power of~$\ell$. The group $\cycU{\ell}{H}$ is the intersection of the
closed subgroup $H$ with the kernel of the continuous map~$\cyc{F}{\Omega}{\ell}{\cdot}$, hence closed in the
compact group~$\GalG{F}{\Omega}$, hence compact. A locally constant map on a compact space takes finitely many
values. So one power of $\ell$ kills all values of a given~$\varphi$.

(b) The conditions for $\Dinf{\ell}{1}{H}$ are those for $\Dinf{\ell}{N}{H}$ without the last one, because $1$
lies in no maximal ideal. An inclusion of modules restricts to the elements killed by~$\ell$, and the length of
a submodule is at most that of the module.

(c) Let \m{\varphi\in\Dinf{\ell}{N}{H}} and \m{u\in\cycU{\ell}{H}}. The elements $g^{-1}$, $u$ and $g$ lie
in~$\cycU{\ell}{H}$, and $\varphi$ is additive, so
\[
\bigl(\dRep{\ell}{N}{H}(g)\,\varphi\bigr)(u)=\varphi(g^{-1}ug)=\varphi(g^{-1})+\varphi(u)+\varphi(g)=\varphi(u),
\]
because \m{\varphi(g^{-1})+\varphi(g)=\varphi(1)=0}. Hence $\dRep{\ell}{N}{H}(g)$ fixes every~$\varphi$. Then
\[
\bigl(\xRep{\ell}{N}{H}(g)\,x\bigr)(\varphi)=x\bigl(\dRep{\ell}{N}{H}(g^{-1})\,\varphi\bigr)=x(\varphi)
\qquad\text{for }x\in\Xinf{\ell}{N}{H},
\]
as $g^{-1}$ lies in~$\cycU{\ell}{H}$ too. Finally \m{\cyc{F}{\Omega}{\ell}{g}=1}, so the twists change nothing.
\end{proof}
\leanline{\lean{dinfTors}; the length in (b) is \lean{xRank}}

In (a), Theorem~A of the second note \cite{N2} proves the same statement; the archive has it in both units.

% ------------------------------------------------------------------------------------- 5
\sect{Generators}{s:gen}
\statement{inksoft}{\phantomsection\label{i:G}\textbf{Input G} (closed subgroups of the units).\enspace\emph{Let
$\ell$ be a prime number and $W$ a closed subgroup of\/~$\Z_\ell^\times$, and assume
\m{W\subseteq\princ{\ell}{2}} if \m{\ell=2}. Then there is \m{s\in W} such that the set of powers \m{s^j},
\m{j\in\Z}, is dense in~$W$.}}
\leanline{\lean{kumExistsGen}}

This is proved in another unit of the archive and is not deformalized here. It expresses the classical
structure of the units: for odd $\ell$ the group $\Z_\ell^\times$ is the product of a cyclic group of order
\m{\ell-1} and of \m{\princ{\ell}{1}\cong\Z_\ell}, and \m{\princ{2}{2}\cong\Z_2} for \m{\ell=2}
\cite[Ch.\,II, §3]{Se}; a closed subgroup of such a group is topologically generated by one element.

\begin{proposition}\label{p:gen}
In Setting~\ref{set:base}, let $H$ be a closed subgroup of~$\GalG{F}{\Omega}$. Assume that
$\ell$ is odd or that \m{H\subseteq\gamU{\ell}{\GalG{F}{\Omega}}}, that is,
\m{\cyc{F}{\Omega}{\ell}{H}\subseteq\princ{\ell}{\qExp{\ell}}}. Then there is an $\ell$-generator for~$H$.
\end{proposition}
\begin{proof}
The group $H$ is closed in the compact group~$\GalG{F}{\Omega}$, so it is compact, and its image
$\cyc{F}{\Omega}{\ell}{H}$ under the continuous map~$\cyc{F}{\Omega}{\ell}{\cdot}$ is compact, hence closed
in~$\Z_\ell^\times$. If \m{\ell=2}, it lies in~$\princ{\ell}{2}$, because \m{\qExp{\ell}=2}. By Input~G there is
\m{s\in\cyc{F}{\Omega}{\ell}{H}} whose powers are dense in~$\cyc{F}{\Omega}{\ell}{H}$. Any \m{\gamma\in H} with
\m{\cyc{F}{\Omega}{\ell}{\gamma}=s} is an $\ell$-generator for~$H$.
\end{proof}
\leanline{\lean{existsGen}}

\begin{proposition}\label{p:genE}
In Setting~\ref{set:ext} the following hold.
\begin{enumerate}[label=\textup{(\alph*)},topsep=3pt,itemsep=1pt]
\item $\galOf{F}{E}{\Omega}$ is open and closed in~$\GalG{F}{\Omega}$, and
\m{\cyc{E}{\Omega}{\ell}{g}=\cyc{F}{\Omega}{\ell}{\galEmb{F}{E}{\Omega}{g}}} for every \m{g\in\GalG{E}{\Omega}}.
\item Condition~$\Procyc{E}{\Omega}{\ell}$ holds if and only if $\ell$ is odd or
\m{\galOf{F}{E}{\Omega}\subseteq\gamU{\ell}{\GalG{F}{\Omega}}}.
\item If condition~$\Procyc{E}{\Omega}{\ell}$ holds, there is an $\ell$-generator for~$\galOf{F}{E}{\Omega}$.
\end{enumerate}
\end{proposition}
\begin{proof}
(a) The group $\galOf{F}{E}{\Omega}$ is one of the basic neighbourhoods of~$1$, so it is open; its complement is a
union of cosets of it, so it is closed. Both characters describe the action on the same roots of unity
of~$\Omega$.
(b) For \m{\ell=2} condition~$\Procyc{E}{\Omega}{\ell}$ says
\m{\cyc{E}{\Omega}{\ell}{\GalG{E}{\Omega}}\subseteq\princ{\ell}{2}}, which by~(a) and Definition~\ref{d:ext} means
\m{\galOf{F}{E}{\Omega}\subseteq\Un{F}{\Omega}{\ell}{2}}, and \m{\qExp{\ell}=2}. (c) This is
Proposition~\ref{p:gen} for \m{H=\galOf{F}{E}{\Omega}}, by (a) and~(b).
\end{proof}
\leanline{\lean{genGal}\sep the subgroup of~$E$: \lean{galOf}, the image of \lean{galEmb}\sep (a) by
\lean{galEmbTop} and \lean{galEmbFam}}

% ------------------------------------------------------------------------------------- 6
\sect{Proof of Theorem B}{s:proofB}
Let $H$, $N$, $h$, $V$ and $\gamma$ be as in the theorem, and let \m{\varphi\in\Dinf{\ell}{N}{H}} with
\m{\dTwOp{\ell}{N}{H}{h}(\gamma)\,\varphi=\varphi}. We show that \m{\dTwOp{\ell}{N}{H}{h}(g)\,\varphi=\varphi} for
every \m{g\in V}.

\emph{A level that fixes~$\varphi$.} The group $H$ is closed and normal, so by Theorem~A(i) of the second note \cite{N2} there is
\m{n\ge0} with \m{\dRep{\ell}{N}{H}(w)\,\varphi=\varphi} for all \m{w\in\levU{\ell}{H}{n}}. By
Proposition~\ref{p:tors}(a) there is \m{r\ge0} with \m{\ell^r\varphi=0}. Put \m{m=\max(n,r)}, and let
\m{w\in\levU{\ell}{H}{m}}. Then \m{\dRep{\ell}{N}{H}(w)\,\varphi=\varphi}, because
\m{\levU{\ell}{H}{m}\subseteq\levU{\ell}{H}{n}}. Moreover \m{\cyc{F}{\Omega}{\ell}{w}^{-h}} lies in
\m{\princ{\ell}{m}\subseteq\princ{\ell}{r}}, so \m{\cyc{F}{\Omega}{\ell}{w}^{-h}=1+\ell^rs} with \m{s\in\Z_\ell}, and
\[
\dTwOp{\ell}{N}{H}{h}(w)\,\varphi=(1+\ell^rs)\,\varphi=\varphi .
\]

\emph{Powers of~$\gamma$.} The \m{g\in\GalG{F}{\Omega}} with \m{\dTwOp{\ell}{N}{H}{h}(g)\,\varphi=\varphi} form a
subgroup of~$\GalG{F}{\Omega}$, because $\dTwOp{\ell}{N}{H}{h}$ is a representation. It contains~$\gamma$, hence
every power~$\gamma^j$.

\emph{Elements of~$V$.} Let \m{g\in V}. The set \m{\cyc{F}{\Omega}{\ell}{g}\,\princ{\ell}{m}} is an open
neighbourhood of~$\cyc{F}{\Omega}{\ell}{g}$ in~$\Z_\ell^\times$, and $\cyc{F}{\Omega}{\ell}{g}$ lies in the closure of
the powers of~$\cyc{F}{\Omega}{\ell}{\gamma}$. So there is an integer~$j$ with
\m{\cyc{F}{\Omega}{\ell}{\gamma}^j\in\cyc{F}{\Omega}{\ell}{g}\,\princ{\ell}{m}}. Put \m{w=\gamma^{-j}g}. Then
\m{w\in V\subseteq H} and \m{\cyc{F}{\Omega}{\ell}{w}\in\princ{\ell}{m}}, so \m{w\in\levU{\ell}{H}{m}}. Since
\m{g=\gamma^jw},
\[
\dTwOp{\ell}{N}{H}{h}(g)\,\varphi=\dTwOp{\ell}{N}{H}{h}(\gamma^j)\,\dTwOp{\ell}{N}{H}{h}(w)\,\varphi
=\dTwOp{\ell}{N}{H}{h}(\gamma^j)\,\varphi=\varphi .
\]

This shows \m{\ker\bigl(\dTwOp{\ell}{N}{H}{h}(\gamma)-1\bigr)\subseteq\ShaD{\ell}{N}{H}{V}{h}}. The other inclusion
holds because \m{\gamma\in V}. Finally $\dTwOp{\ell}{N}{H}{h}(\gamma^{-1})$ is the inverse
of~$\dTwOp{\ell}{N}{H}{h}(\gamma)$, so the two operators fix the same elements. \qed
\leanline{\lean{shaDGen}, with \lean{dinfNil} for the integer~$n$}

% ------------------------------------------------------------------------------------- 7
\sect{Transport of characters}{s:transport}
The archive compares the modules of characters of two number fields with two algebraic closures.

\begin{setting}\label{set:tr}
$F$ and $F'$ are number fields, $\Omega$ is an algebraic closure of~$F$ and $\Omega'$ is an algebraic closure
of~$F'$, \m{f\colon F\to F'} is a homomorphism of fields, and \m{\tau\colon\Omega\to\Omega'} is an isomorphism of
fields that extends~$f$: \m{\tau(t)=f(t)} for all \m{t\in F}.
\end{setting}

\begin{lemma}\label{l:theta}
In Setting~\ref{set:tr}, for \m{g'\in\GalG{F'}{\Omega'}} the map \m{\theta(g')=\tau^{-1}g'\tau} is an automorphism
of~$\Omega$ that fixes every element of~$F$. The map
\[
\theta\colon\GalG{F'}{\Omega'}\to\GalG{F}{\Omega},\qquad \tau\bigl(\theta(g')\,t\bigr)=g'(\tau t)\quad\text{for }t\in\Omega ,
\]
is an injective continuous homomorphism, and it is the only map with this property. For every prime number~$\ell$
and every \m{g'\in\GalG{F'}{\Omega'}} we have \m{\cyc{F}{\Omega}{\ell}{\theta(g')}=\cyc{F'}{\Omega'}{\ell}{g'}}.
\end{lemma}
\begin{proof}
The map \m{\tau^{-1}g'\tau} fixes~$F$, because $g'$ fixes~$f(F)$. The displayed equation determines
$\theta(g')$, and $\theta$ is a homomorphism; it is injective because $\tau$ is bijective. It is continuous: the
elements that $\theta$ maps into~$\galOf{F}{E}{\Omega}$, for a finite extension $E$ of $F$ in~$\Omega$, are those that
fix~$\tau(E)$, and these form an open subgroup of~$\GalG{F'}{\Omega'}$. Since $\tau$ maps
roots of unity to roots of unity, \m{\cyc{F}{\Omega}{\ell}{\theta(g')}} and \m{\cyc{F'}{\Omega'}{\ell}{g'}} are the
same exponent.
\end{proof}
\leanline{\lean{galTransport}; the formal statement has one more clause, not used here: $\theta$ is surjective if
$f$ is}

\begin{lemma}\label{l:val}
In Setting~\ref{set:tr}, let $\theta$ be the homomorphism of Lemma~\ref{l:theta}, let $A'$ be a valuation ring
of~$\Omega'$, and let \m{A=\tau^{-1}(A')}, a valuation ring of~$\Omega$. The groups $\Dec{A}$ and $\Ine{A}$ are
formed in~$\GalG{F}{\Omega}$, the groups $\Dec{A'}$ and $\Ine{A'}$ in~$\GalG{F'}{\Omega'}$.
\begin{enumerate}[label=\textup{(\alph*)},topsep=3pt,itemsep=1pt]
\item An element \m{t\in\Omega} lies in $\mathfrak m_A$ if and only if $\tau(t)$ lies in~$\mathfrak m_{A'}$; that is,
\m{\mathfrak m_A=\tau^{-1}(\mathfrak m_{A'})}. In particular an integer lies in $\mathfrak m_A$ if and only if it
lies in~$\mathfrak m_{A'}$.
\item Let \m{g'\in\GalG{F'}{\Omega'}} and \m{g=\theta(g')}, so that \m{\tau(g(t))=g'(\tau(t))} for \m{t\in\Omega}.
Then \m{g\in\Dec{A}} if and only if \m{g'\in\Dec{A'}}, and \m{g\in\Ine{A}} if and only if \m{g'\in\Ine{A'}}.
\end{enumerate}
\end{lemma}
\begin{proof}
(a) An element $t$ lies in $\mathfrak m_A$ if and only if \m{t\in A} and either \m{t=0} or \m{t^{-1}\notin A}. The
isomorphism $\tau$ respects zero and inverses and maps $A$ onto~$A'$. It fixes the integers.
(b) Since \m{g=\tau^{-1}g'\tau}, we have \m{g(A)=\tau^{-1}\bigl(g'(A')\bigr)}, which equals \m{A=\tau^{-1}(A')} if and
only if \m{g'(A')=A'}. For \m{t\in A} we have \m{\tau(g(t)-t)=g'(\tau t)-\tau t}, and $\tau t$ runs through~$A'$. By (a)
the left side lies in $\mathfrak m_{A'}$ exactly when \m{g(t)-t\in\mathfrak m_A}.
\end{proof}
\leanline{(a) \lean{dtVal}\sep (b) \lean{dtDec}, \lean{dtIne}}

\begin{proposition}\label{p:transport}
In Setting~\ref{set:tr}, let \m{\theta\colon\GalG{F'}{\Omega'}\to\GalG{F}{\Omega}} be the homomorphism of
Lemma~\ref{l:theta}, the one with \m{\tau(\theta(g')\,t)=g'(\tau(t))} for all $g'$ and all \m{t\in\Omega}. Let $\ell$
be a prime number and \m{N\ge0} an integer, and let $H$ be a subgroup of~$\GalG{F}{\Omega}$ and $H'$ a subgroup
of~$\GalG{F'}{\Omega'}$ with \m{\theta(\cycU{\ell}{H'})=\cycU{\ell}{H}}. Then the map
\m{\varphi\mapsto\varphi\circ\theta} is an isomorphism of\/ $\Z_\ell$-modules from $\Dinf{\ell}{N}{H}$
onto~$\Dinf{\ell}{N}{H'}$.
\end{proposition}
\begin{proof}
A continuous injective map from a compact space to a Hausdorff space is a homeomorphism onto its image. So
$\theta$ restricts to an isomorphism of topological groups from $\cycU{\ell}{H'}$ onto~$\cycU{\ell}{H}$. Let
\m{\varphi\in\Dinf{\ell}{N}{H}}. Then \m{\varphi\circ\theta} is locally constant and additive. Let $A'$ be a
valuation ring of~$\Omega'$ and \m{A=\tau^{-1}(A')}, with the four groups formed as in Lemma~\ref{l:val}. If
\m{g'\in\cycU{\ell}{H'}} lies in~$\Ine{A'}$, then $\theta(g')$ lies in \m{\cycU{\ell}{H}\cap\Ine{A}} by
Lemma~\ref{l:val}(b), so \m{\varphi(\theta(g'))=0}. If \m{N\in\mathfrak m_{A'}} and $g'$ lies in~$\Dec{A'}$, then
\m{N\in\mathfrak m_A} and \m{\theta(g')\in\cycU{\ell}{H}\cap\Dec{A}}, so again \m{\varphi(\theta(g'))=0}. Hence
\m{\varphi\circ\theta} lies in~$\Dinf{\ell}{N}{H'}$. In the same way composition with the inverse of the
restricted map carries the second module into the first: every valuation ring of $\Omega$ is
\m{\tau^{-1}(A')} for the valuation ring \m{A'=\tau(A)}, and Lemma~\ref{l:val} is an equivalence. The two maps are
linear and inverse to each other.
\end{proof}
\leanline{\lean{dinfTransport}, with \lean{dtMem} for the four conditions}

We use the proposition in one case. In Setting~\ref{set:ext}, take $E$ for the second field, $\Omega$ for both
algebraic closures, the inclusion of $F$ in~$E$ for the homomorphism of fields and the identity of~$\Omega$ for
the isomorphism. Then the map of Lemma~\ref{l:theta} is \m{\galEmb{F}{E}{\Omega}{\cdot}\colon\GalG{E}{\Omega}\to\GalG{F}{\Omega}},
and the proposition says that for subgroups $H$ of~$\GalG{F}{\Omega}$ and $H'$ of~$\GalG{E}{\Omega}$ such that
this map carries $\cycU{\ell}{H'}$ onto~$\cycU{\ell}{H}$, composition with it identifies $\Dinf{\ell}{N}{H}$
with~$\Dinf{\ell}{N}{H'}$: it does not matter whether the decomposition and inertia groups are taken
in~$\GalG{E}{\Omega}$ or in~$\GalG{F}{\Omega}$. This case can also be seen directly, since an element $g$
of~$\GalG{E}{\Omega}$ lies in the group $\Dec{A}$ or $\Ine{A}$ formed in~$\GalG{E}{\Omega}$ exactly when
$\galEmb{F}{E}{\Omega}{g}$ lies in the group $\Dec{A}$ or $\Ine{A}$ formed in~$\GalG{F}{\Omega}$.

% ------------------------------------------------------------------------------------- 8
\sect{Proof of Theorem A}{s:proofA}
\begin{definition}[the classes]\label{d:classes}
In Setting~\ref{set:base}, let $S$ be the set of the prime ideals of the ring of integers of~$F$ that
contain~$\ell$. Let $\GS{\Omega}{S}$ be the quotient of $\GalG{F}{\Omega}$ by the closed normal subgroup generated by
the inertia groups $\Ine{A}$, for the valuation rings \m{A\ne\Omega} of~$\Omega$ with \m{\ell\notin\mathfrak m_A}; it
is the Galois group of the largest extension of $F$ in $\Omega$ that is unramified at the primes outside~$S$. The
cyclotomic character $\cyc{F}{\Omega}{\ell}{\cdot}$ factors through~$\GS{\Omega}{S}$.

A \kw{datum} for $\ell$, $F$ and~$\Omega$ consists of the set~$S$, of a valuation ring $A_{\mathfrak p}$ of~$\Omega$
for each \m{\mathfrak p\in S} whose maximal ideal contains~$\mathfrak p$, and of the homomorphism
\m{\chi\colon \GS{\Omega}{S}\to\Z_\ell^\times} whose composite with the projection
\m{\GalG{F}{\Omega}\to \GS{\Omega}{S}} is~$\cyc{F}{\Omega}{\ell}{\cdot}$; the archive carries this factorization as
part of its data. The set of these data is~$\LData{\ell}{F}{\Omega}$. Let \m{\mathcal D\in\LData{\ell}{F}{\Omega}} be a
datum, with valuation rings $A_{\mathfrak p}$ and homomorphism~$\chi$, and let $h$ be an integer.
\begin{enumerate}[label=\textup{(\roman*)},topsep=3pt,itemsep=1pt]
\item $\qzTw{\mathcal D}{h}$ is the module $\QZl{\ell}$ on which \m{g\in \GS{\Omega}{S}} acts as multiplication
by~\m{\chi(g)^{-h}}. The group \m{H^1(\qzTw{\mathcal D}{h})} consists of the continuous maps $z$ from
$\GS{\Omega}{S}$ to~$\QZl{\ell}$ with \m{z(gw)=z(g)+\chi(g)^{-h}z(w)}, modulo those of the form
\m{g\mapsto\chi(g)^{-h}y-y} with \m{y\in\QZl{\ell}}.
\item $\ShaOf{\mathcal D}{h}$ is the submodule of the classes in \m{H^1(\qzTw{\mathcal D}{h})} whose restriction to
the image of the decomposition group $\Dec{A_{\mathfrak p}}$ in~$\GS{\Omega}{S}$ is trivial for each
\m{\mathfrak p\in S}
\cite[§7]{N1}. These are the \kw{locally trivial classes}.
\item On the kernel $\cycU{\ell}{\GalG{F}{\Omega}}$ of~$\cyc{F}{\Omega}{\ell}{\cdot}$ the action on
$\qzTw{\mathcal D}{h}$ is trivial. So the composite of a map $z$ as in~(i) with the projection
\m{\GalG{F}{\Omega}\to \GS{\Omega}{S}}, restricted to~$\cycU{\ell}{\GalG{F}{\Omega}}$, is a homomorphism
\m{\cycU{\ell}{\GalG{F}{\Omega}}\to\QZl{\ell}}, and it depends only on the class $\xi$ of~$z$. We call
it~$\resInf{\mathcal D}{h}{\xi}$.
\item $\qzG{F}{\Omega}{\ell}{h}$ is the module $\QZl{\ell}$ on which \m{g\in\GalG{F}{\Omega}} acts as
multiplication by~\m{\cyc{F}{\Omega}{\ell}{g}^{-h}}.
\end{enumerate}
\end{definition}
\leanline{\lean{LData}\sep\lean{qzTw}\sep\lean{Sha}\sep\lean{resInf}\sep the action in~(iv): \lean{qzG}}

\begin{definition}[equivariant characters]\label{d:eqhom}
In Setting~\ref{set:base}, let $C$ be a subgroup of~$\GalG{F}{\Omega}$, let \m{N\ge0} be an integer, $\Lambda$ a
ring, $M$ a $\Lambda$-module and $\pi$ a representation of $\GalG{F}{\Omega}$ on~$M$. Then $\EqHom{N}{C}{\pi}$ is
the $\Lambda$-module of all \m{\psi\in\XDual{N}{C}{\Lambda}{M}} such that
\[
\psi(gug^{-1})=\pi(g)\,\psi(u)\qquad\text{for all }g\in\GalG{F}{\Omega}\text{ and }u\in C\text{ with }gug^{-1}\in C .
\]
\end{definition}
\leanline{\lean{EqHom}}

\statement{inksoft}{\phantomsection\label{i:R1}\textbf{Input R1} (restriction is injective).\enspace\emph{In
Setting~\ref{set:base}, let \m{\mathcal D\in\LData{\ell}{F}{\Omega}} be a datum (Definition~\ref{d:classes}) and
\m{h\ne0} an integer, and assume that condition~$\Procyc{F}{\Omega}{\ell}$ holds. Then
\m{\xi\mapsto\resInf{\mathcal D}{h}{\xi}} is injective on \m{H^1(\qzTw{\mathcal D}{h})}.}

\vspace{3pt}\phantomsection\label{i:R2}\textbf{Input R2} (the image of restriction).\enspace\emph{In
Setting~\ref{set:base}, let \m{\mathcal D\in\LData{\ell}{F}{\Omega}} be a datum and \m{h\ne0} an integer, and
assume that condition~$\Procyc{F}{\Omega}{\ell}$ holds. Then the image of\/ $\ShaOf{\mathcal D}{h}$ under
\m{\xi\mapsto\resInf{\mathcal D}{h}{\xi}} is the module
$\EqHom{\ell}{\cycU{\ell}{\GalG{F}{\Omega}}}{\qzG{F}{\Omega}{\ell}{h}}$ of Definition~\ref{d:eqhom}: the module of
all \m{\psi\in\Dinf{\ell}{\ell}{\GalG{F}{\Omega}}} such that}
\[
\psi(gug^{-1})=\cyc{F}{\Omega}{\ell}{g}^{-h}\,\psi(u)\qquad\text{\emph{for all }}g\in\GalG{F}{\Omega},\ u\in\cycU{\ell}{\GalG{F}{\Omega}} .
\]}
\leanline{R1 \lean{resInfInj}\sep R2 \lean{resInfOnto}, with the module \lean{EqHom}}

Both are theorems of the unit of the archive on Kummer theory, the unit of \cite{N2}; they are not deformalized
here. The classical reason is the inflation and restriction sequence for the
subgroup~$\cycU{\ell}{\GalG{F}{\Omega}}$: the quotient of $\GalG{F}{\Omega}$ by it is isomorphic to the group
$\cyc{F}{\Omega}{\ell}{\GalG{F}{\Omega}}$, which is topologically generated by one element under
condition~$\Procyc{F}{\Omega}{\ell}$, and for \m{h\ne0} this element acts on $\qzG{F}{\Omega}{\ell}{h}$ by a scalar
different from~$1$ \cite[Ch.\,I, §6]{NSW}.

\begin{definition}[given by restriction]\label{d:desc}
In Setting~\ref{set:ext}, let \m{\mathcal D\in\LData{\ell}{E}{\Omega}} be a datum for the field~$E$, $h$ an integer,
$H$ a normal subgroup of~$\GalG{F}{\Omega}$, and
\[
\Phi\colon\ \ShaOf{\mathcal D}{h}\ \longrightarrow\ \ShaD{\ell}{\ell}{H}{\galOf{F}{E}{\Omega}}{h}
\]
a homomorphism of $\Z_\ell$-modules. We say that $\Phi$ is \kw{given by restriction for $E$, $\mathcal D$, $h$
and~$H$} if \m{(\Phi\xi)(u')=\resInf{\mathcal D}{h}{\xi}(u)} for all \m{\xi\in\ShaOf{\mathcal D}{h}}, all
\m{u\in\cycU{\ell}{\GalG{E}{\Omega}}} and all \m{u'\in\cycU{\ell}{H}} with \m{u'=\galEmb{F}{E}{\Omega}{u}}.
\end{definition}
\leanline{\lean{IsDesc}}

\begin{proof}[Proof of Theorem A]
Inputs R1 and~R2 are used for the number field~$E$ with its algebraic closure~$\Omega$, for the
datum~$\mathcal D$ and the integer~$h$; their hypotheses are among those of the theorem. By
Proposition~\ref{p:genE}(a) and Definition~\ref{d:ext}, the map $\galEmb{F}{E}{\Omega}{\cdot}$ carries
$\cycU{\ell}{\GalG{E}{\Omega}}$ bijectively onto~$\cycU{\ell}{\galOf{F}{E}{\Omega}}$, and this is $\cycU{\ell}{H}$ by
hypothesis. By Proposition~\ref{p:transport}, in the case described after it, composition with this map is an
isomorphism from \m{\Dinf{\ell}{\ell}{H}} onto the module $\Dinf{\ell}{\ell}{\GalG{E}{\Omega}}$ that is formed
in~$\GalG{E}{\Omega}$.

Let \m{\varphi\in\Dinf{\ell}{\ell}{H}} and \m{g\in\galOf{F}{E}{\Omega}}. The equation
\m{\dTwOp{\ell}{\ell}{H}{h}(g)\,\varphi=\varphi} says that
\[
\cyc{F}{\Omega}{\ell}{g}^{-h}\,\varphi(g^{-1}ug)=\varphi(u)\qquad\text{for all }u\in\cycU{\ell}{H} .
\]
The group $\cycU{\ell}{H}$ is normal in~$\GalG{F}{\Omega}$, so we may write \m{u=gwg^{-1}} with
\m{w\in\cycU{\ell}{H}}, and the equation becomes
\[
\varphi(gwg^{-1})=\cyc{F}{\Omega}{\ell}{g}^{-h}\,\varphi(w)\qquad\text{for all }w\in\cycU{\ell}{H} .
\]
So the isomorphism above carries the module of invariants onto the module of Input~R2 for the field~$E$:
\[
\ShaD{\ell}{\ell}{H}{\galOf{F}{E}{\Omega}}{h}\ \longrightarrow\
\EqHom{\ell}{\cycU{\ell}{\GalG{E}{\Omega}}}{\qzG{E}{\Omega}{\ell}{h}} .
\]
By Input~R2 the map \m{\xi\mapsto\resInf{\mathcal D}{h}{\xi}} sends $\ShaOf{\mathcal D}{h}$ onto the second module,
and by Input~R1 it is injective. It is $\Z_\ell$-linear. Composed with the inverse of the isomorphism above it is the map~$\Phi$ of the
theorem, and it is given by restriction for $E$, $\mathcal D$, $h$ and~$H$ by construction.
\end{proof}
\leanline{\lean{desc}}

% ------------------------------------------------------------------------------------- 9
\sect{A separation lemma}{s:sep}
The unit contains one more statement, used in the lemma on characters of the first note \cite[§8]{N1} to show that
the characters of a
$\Z_\ell$-module separate its elements. For a prime number~$\ell$ we write $\vl{\ell}{s}$ for the $\ell$-adic valuation
of an element $s$ of~$\Z_\ell$.

\begin{lemma}\label{l:sep}
Let $\ell$ be a prime number, $M$ a $\Z_\ell$-module and \m{b\in M} a nonzero element. Then there is a nonzero
\m{y\in\QZl{\ell}} such that every \m{s\in\Z_\ell} with \m{sb=0} satisfies \m{sy=0}.
\end{lemma}
\begin{proof}
A nonzero \m{s\in\Z_\ell} is the product of $\ell^{\vl{\ell}{s}}$ and a unit, so it kills an element of a module if and
only if $\ell^{\vl{\ell}{s}}$ does. Suppose first that some power of $\ell$ kills~$b$, and let $\ell^r$ be the smallest
one. Then \m{r\ge1}, because \m{b\ne0}. Let $y$ be the class of $\ell^{-r}$ in~$\QZl{\ell}$; it is nonzero and
\m{\ell^ry=0}. If \m{s\ne0} and \m{sb=0}, then \m{\ell^{\vl{\ell}{s}}b=0}, so \m{\vl{\ell}{s}\ge r}, so
\m{\ell^{\vl{\ell}{s}}y=0} and \m{sy=0}. If no power of $\ell$ kills~$b$, then only \m{s=0} kills~$b$, and the class of
$\ell^{-1}$ will do.
\end{proof}
\leanline{\lean{qzAnn}; the classes of $\ell^{-r}$ come from \lean{zmodToQZ}}

In \cite{N1} the lemma is applied as follows: \m{sb\mapsto sy} is a well-defined homomorphism from the submodule
generated by $b$ to~$\QZl{\ell}$, and it extends to~$M$ because $\QZl{\ell}$ is a divisible group.

% ------------------------------------------------------------------------------------- 10
\sect{Where the theorems are used}{s:use}
In \cite[§8]{N1} the base field is called $L$ and the intermediate field~$F$; they are the fields $F$ and $E$ of
this note. Input~2a there is Theorem~A for \m{H=\GalG{F}{\Omega}}, and Input~2b is Theorem~B for
\m{H=\GalG{F}{\Omega}}, with the element $\eta$ there in the place of~$\gamma$. With them the first note computes
the length of the group $\ShaOf{\mathcal D}{h}$ of a datum~$\mathcal D$ for the layer \m{E=\layF{\ell}{F}{\Omega}{m}} of
index~$m$ as the length of the kernel of \m{\dTwOp{\ell}{\ell}{\GalG{F}{\Omega}}{h}(\eta)-1}
on~$\dTw{\ell}{\ell}{\GalG{F}{\Omega}}{h}$, where $\eta$ is a power of an $\ell$-generator for~$\GalG{F}{\Omega}$, and by duality as the length
of a cokernel on~$\xTw{\ell}{\ell}{\GalG{F}{\Omega}}{h}$. Proposition~\ref{p:roots} and Proposition~\ref{p:genE} enter
the description of the layers there, the lemmas on roots of unity and on the layers \cite[§8]{N1}, and the tower
condition of Definition~\ref{d:tower} is the
hypothesis of its main theorem. With this note, of the three inputs of \cite{N1} the first is proved
in~\cite{N2}; the second is reduced to Inputs G, R1 and~R2 and to the comparison of two descriptions of the
length of the group of locally trivial classes that is cited there; the third, on global duality, has not been
read yet.

% ------------------------------------------------------------------------------------- 11
\sect{Notes on the formal text}{s:notes}
\begin{enumerate}[label=\textup{\arabic*.},leftmargin=1.6em,itemsep=5pt]
\item \emph{Theorem A is an assembly.} Its formal proof transports the characters from $\GalG{F}{\Omega}$
to~$\GalG{E}{\Omega}$, rewrites the invariance as the equation of Input~R2, and applies the two theorems on
restriction. The hypotheses \m{h\ne0} and condition~$\Procyc{E}{\Omega}{\ell}$ are used in those two theorems only.
\item \emph{The form of Theorem A.} The theorem is stated as in the archive: there is an isomorphism $\Phi$ from
$\ShaOf{\mathcal D}{h}$ to the module of invariants such that \m{(\Phi\xi)(u')} is the value of
$\resInf{\mathcal D}{h}{\xi}$ at~$u$, whenever \m{u'\in\cycU{\ell}{H}} is the element $u$
of~$\cycU{\ell}{\GalG{E}{\Omega}}$ regarded in~$\GalG{F}{\Omega}$. Every element of $\cycU{\ell}{H}$ arises so from
exactly one~$u$, hence $\Phi$ is the restriction map, as is said after the theorem.
\item \emph{Two algebraic closures.} Proposition~\ref{p:transport} is stated for two fields with two algebraic
closures and an isomorphism between them. The archive uses it once, with the identity.
\item \emph{Where the hypotheses of Theorem B enter.} Closedness and normality of $H$ are used for the integers
$n$ and~$r$. The group $V$ need not be closed or normal, and $\gamma$ need not generate it in any other sense
than Definition~\ref{d:gen}.
\item \emph{The integer $\cZero{\ell}{F}$} is defined in the archive as the supremum of a set of natural numbers,
which is meaningful because of Lemma~\ref{l:deg}. The definition is made there for every field and every natural
number in the place of the prime.
\item \emph{Lemma~\ref{l:val}} is proved in the archive in one direction, and the other direction is obtained
where it is needed by exchanging the two fields. The equivalence is stated here at once.
\item \emph{Facts from other units.} The properties of the map~$\theta$ in Lemma~\ref{l:theta} and of the group
$\galOf{F}{E}{\Omega}$ in Proposition~\ref{p:genE}(a) are theorems of other units of the archive. Their formal proofs
were not read; the short arguments given here are the standard ones.
\item \emph{The tower condition} is stated in the archive for the module over the algebraic closure of~$F$ that
the library constructs, not over a closure that is given. The modules over two algebraic closures of~$F$ are
isomorphic (Proposition~\ref{p:transport}, for an isomorphism of the two closures over~$F$).
\item \emph{The unit.} The \UnitMembers{} declarations treated here are those of the unit that the final
theorems use. The files of the unit contain further declarations, on norms of ideal classes, on coinvariants and
on eigenspaces, that the final theorems do not use; for a few of them the reading project could not decide this.
\end{enumerate}

\subhead{How this note was written.} The reading project supplied a dossier of the unit: the statements, a
glossary, readings of the imported declarations, and a structured proof of every theorem in which each step
quotes the lines of the archive that it renders. A program compared the quotations with the archive, for all
\UnitSteps{} steps of the \UnitProofs{} proofs. The proofs of this note were written from those structured proofs.
The formal statements of the theorems and the file of definitions were read in the archive and compared with the
statements above; of the formal proofs, the opening lines of the proof of Theorem~B were read, to identify the two
results it cites. The arguments for Lemma~\ref{l:fix} and Lemma~\ref{l:val} are shorter here than in the
archive and are ours.

\subhead{Changes in this edition.} The earlier version of this note fixed the prime, the number field and its
algebraic closure once and then wrote the modules, the groups, the characters and the operators without them: a
bare letter with the subscript~$\infty$ for the modules $\Dinf{\ell}{N}{H}$ and~$\Xinf{\ell}{N}{H}$, a subscript for
the levels $\levU{\ell}{H}{n}$, one letter each for the cyclotomic character, the exponent~$\cZero{\ell}{F}$ and the
operators $\dTwOp{\ell}{N}{H}{h}(g)$, and the group of locally trivial classes without its datum. Definitions stood
in running text, and hypotheses in sentences before the statements. Now every function that the archive defines
is written with all the arguments of its formal definition, the definitions are numbered, and every statement
states its hypotheses or cites the Setting that holds them. Three things changed beyond notation. Theorem~A is
stated in the form of the archive, with the isomorphism~$\Phi$ and Definition~\ref{d:desc}; its earlier form,
that restriction is an isomorphism onto the invariants, is the paragraph after it. The cohomology classes and
Inputs R1 and~R2 are stated for the field of Setting~\ref{set:base}, as in the archive, and applied to~$E$ in the
proof of Theorem~A. The tower condition (Definition~\ref{d:tower}) names the algebraic closure over which its
module is formed, which the earlier text left out. No proof changed in content.

% ------------------------------------------------------------------------------------- A
\head{Appendix A. Concordance with the Lean archive}
{\fontsize{9.6}{12.8}\selectfont\color{inksoft}\noindent Each name is a link to the lines of the declaration at
commit {\leanfont\Pin} of the public mirror of the archive. The line ranges and the quotations below are read from
the archive by the build script, which also checks them against the SHA-256 of each file. The first
\UnitMembers{} rows are the declarations of the unit.\par}
\vspace{4pt}
{\fontsize{8.6}{11.6}\selectfont\renewcommand{\leanfont}{\ttfamily\fontsize{7.7}{10}\selectfont\color{leancol}}
\setlength{\LTleft}{0pt}\setlength{\LTright}{\fill}\setlength{\tabcolsep}{0pt}
\begin{longtable}{@{}p{30mm}p{56mm}p{43mm}r@{}}
\textbf{In this note} & \textbf{Lean declaration} & \textbf{File} & \textbf{Lines}\\\midrule\endhead
\ConcordanceRows
\end{longtable}}

\Needspace{8\baselineskip}
\subhead{Three statements as they are written in the archive.} Theorem~A; here {\leanfont k} and {\leanfont L}
are the fields $F$ and~$E$, {\leanfont U} is the group~$H$, {\leanfont galOf k L Ω}
is~$\galOf{F}{E}{\Omega}$, {\leanfont e} is~$h$, {\leanfont D} is the datum~$\mathcal D$, and
{\leanfont Procyc L Ω l} is condition~$\Procyc{E}{\Omega}{\ell}$:
\quoted{desc}
Theorem~B; here {\leanfont dTw l N U e γ} is the operator~$\dTwOp{\ell}{N}{H}{h}(\gamma)$:
\quoted{shaDGen}
Proposition~\ref{p:roots}; here {\leanfont cZero l L} is the integer~$\cZero{\ell}{F}$ of the field, which the
archive calls~{\leanfont L}, {\leanfont Un L Ω l c} is the group~$\Un{F}{\Omega}{\ell}{c}$ and {\leanfont qExp l}
is~$\qExp{\ell}$:
\quoted{muQ}

% ------------------------------------------------------------------------------------- notation
\head{Notation}
{\fontsize{9.6}{12.6}\selectfont\color{inksoft}\noindent The symbols of the text. In brackets: the function of
the archive with its arguments, or the name that the archive uses.\par}
\vspace{4pt}
\begingroup\fontsize{8.6}{10.9}\selectfont\raggedright
\renewcommand{\leanfont}{\ttfamily\fontsize{7.6}{10}\selectfont\color{leancol}}
\setlength{\columnsep}{9mm}\setlength{\parskip}{0pt}
\begin{multicols}{2}
\nt{$\ell$, $\Z_\ell$, $\Q_\ell$}{a prime [{\leanfont l}]; the $\ell$-adic integers and numbers}
\nt{$\QZl{\ell}$}{the quotient group of $\Q_\ell$ by~$\Z_\ell$ [{\leanfont QZ l}]}
\nt{$F$, $\Omega$, $\GalG{F}{\Omega}$}{a number field [{\leanfont k}, {\leanfont F}, {\leanfont L}]; an algebraic closure; the automorphisms of $\Omega$ over~$F$ [{\leanfont G F Ω}]}
\nt{$\cyc{F}{\Omega}{\ell}{g}$}{the $\ell$-adic cyclotomic character of $\Omega$ over~$F$ at~$g$ [{\leanfont cycZ F Ω l g}]}
\nt{$\princ{\ell}{n}$, $W$}{the units congruent to $1$ modulo~$\ell^n$ [{\leanfont princ l n}]; a closed subgroup of~$\Z_\ell^\times$}
\nt{$H$, $\levU{\ell}{H}{n}$, $\cycU{\ell}{H}$}{a subgroup of $\GalG{F}{\Omega}$ [{\leanfont U}]; its elements with cyclotomic value in~$\princ{\ell}{n}$ [{\leanfont levU l U n}]; those with value~$1$ [{\leanfont cycU l U}]}
\nt{$\Un{F}{\Omega}{\ell}{n}$}{the same for the whole group [{\leanfont Un F Ω l n}]}
\nt{$\qExp{\ell}$, $\gamU{\ell}{H}$}{$2$ for \m{\ell=2}, else $1$ [{\leanfont qExp l}]; the subgroup of $H$ of that level [{\leanfont gamU l U}]}
\nt{$\Procyc{F}{\Omega}{\ell}$}{the condition of Definition~\ref{d:groups}(iv) [{\leanfont Procyc F Ω l}]}
\nt{$E$, $\galOf{F}{E}{\Omega}$}{a finite extension of $F$ in~$\Omega$; the elements of $\GalG{F}{\Omega}$ that fix~$E$ [{\leanfont galOf k E Ω}]}
\nt{$\galEmb{F}{E}{\Omega}{g}$}{an element $g$ of $\GalG{E}{\Omega}$ regarded in~$\GalG{F}{\Omega}$ [{\leanfont galEmb k E Ω g}]}
\nt{$\zeta$, $t$, $\mu_{\ell^n}(\Omega)$}{a root of unity; an element of~$\Omega$; the $t$ with \m{t^{\ell^n}=1}}
\nt{$A$, $\mathfrak m_A$}{a valuation ring of~$\Omega$; its maximal ideal}
\nt{$\Dec{A}$, $\Ine{A}$}{decomposition and inertia group of $A$ in~$\GalG{F}{\Omega}$ [{\leanfont Dec A}, {\leanfont Ine A}]}
\nt{$C$, $N$, $\Lambda$, $M$}{a subgroup of~$\GalG{F}{\Omega}$; an integer \m{\ge0}; a ring; a module over it}
\nt{$\XDual{N}{C}{\Lambda}{M}$}{the characters \m{C\to M} with the conditions of Definition~\ref{d:char} [{\leanfont XDual N U R M}]}
\nt{$\Dinf{\ell}{N}{H}$, $\Xinf{\ell}{N}{H}$}{\m{\XDual{N}{\cycU{\ell}{H}}{\Z_\ell}{\QZl{\ell}}} [{\leanfont Dinf l N U}]; its dual [{\leanfont Xinf l N U}]}
\nt{$\XN{\ell}{F}{\Omega}{N}$}{the dual for the whole group [{\leanfont XN l L Ω N}]}
\nt{$\overline F$, $\XX{\ell}{F}$, $\XP{\ell}{F}$}{the algebraic closure of the library; the duals over it with \m{N=1} [{\leanfont XX l L}] and \m{N=\ell} [{\leanfont XP l L}]}
\nt{$\varphi$, $\psi$, $x$}{characters; an element of~$\Xinf{\ell}{N}{H}$}
\nt{$g$, $\gamma$, $u$, $u'$, $w$}{elements of groups; $\gamma$ an $\ell$-generator (Definition~\ref{d:gen}) [{\leanfont IsGen l U γ}]}
\nt{$\Gamma$, $\pi$, $\omega$, $\twist{\pi}{\omega}$}{a group; a representation; a homomorphism to the units of the ring; the twist of $\pi$ by~$\omega$ [{\leanfont twist π ψ}]}
\nt{$\funRep{\Lambda}{M}{C}$}{conjugation on the maps \m{C\to M} [{\leanfont funRep R M U}]}
\nt{$\dRep{\ell}{N}{H}$, $\xRep{\ell}{N}{H}$}{conjugation on $\Dinf{\ell}{N}{H}$ [{\leanfont dRep l N U}]; the dual action [{\leanfont xRep l N U}]}
\nt{$h$, $\dTwOp{\ell}{N}{H}{h}$, $\xTwOp{\ell}{N}{H}{h}$}{an integer [{\leanfont e}]; the twisted actions (Definition~\ref{d:twist}) [{\leanfont dTw l N U e}, {\leanfont xTw l N U e}]}
\nt{$\dTw{\ell}{N}{H}{h}$, $\xTw{\ell}{N}{H}{h}$}{the two modules with these actions}
\nt{$V$, $\ShaD{\ell}{N}{H}{V}{h}$}{a subgroup of~$\GalG{F}{\Omega}$; the elements fixed by all $\dTwOp{\ell}{N}{H}{h}(g)$, \m{g\in V} [{\leanfont ShaD l N U V e}]}
\nt{$\cZero{\ell}{F}$, $a$}{the largest $c$ with a root of unity of order $\ell^c$ in~$F$ [{\leanfont cZero l L}]; this number in a proof}
\nt{$\layF{\ell}{F}{\Omega}{m}$}{the layer of index~$m$ [{\leanfont layF l L Ω m}]}
\nt{$\layU{\ell}{F}{\Omega}{m}$}{its group [{\leanfont layU l L Ω m}]}
\nt{$\len M$, $\alpha$}{the length of a module; a linear map \m{M\to M}}
\nt{$\xRank{\ell}{N}{H}$}{the length of the elements killed by~$\ell$ [{\leanfont xRank l N U}]}
\nt{$R$}{the number of generators in the tower condition}
\nt{$n$, $m$, $r$, $c$, $j$}{integers}
\nt{$s$, $\vl{\ell}{s}$}{an element of $\Z_\ell$ or of~$\Q_\ell$; its $\ell$-adic valuation [{\leanfont vl l x}]}
\nt{$F'$, $\Omega'$, $H'$, $A'$, $g'$}{a second number field with its algebraic closure, a subgroup, a valuation ring, an element}
\nt{$f$, $\tau$, $\theta$}{a homomorphism \m{F\to F'}; an isomorphism \m{\Omega\to\Omega'}; the map \m{g'\mapsto\tau^{-1}g'\tau}}
\nt{$S$, $\mathfrak p$, $\GS{\Omega}{S}$}{the primes of $F$ above~$\ell$; one of them; the Galois group of the largest extension unramified outside~$S$}
\nt{$\LData{\ell}{F}{\Omega}$, $\mathcal D$}{the data of Definition~\ref{d:classes} [{\leanfont LData F Ω l}]; one of them, with its valuation rings $A_{\mathfrak p}$ and its homomorphism~$\chi$}
\nt{$\qzTw{\mathcal D}{h}$}{$\QZl{\ell}$ with the action through~$\chi^{-h}$ [{\leanfont qzTw D e}]}
\nt{$H^1(\qzTw{\mathcal D}{h})$}{its first cohomology group}
\nt{$z$, $\xi$, $y$}{a crossed homomorphism; its class; an element of~$\QZl{\ell}$}
\nt{$\ShaOf{\mathcal D}{h}$}{the locally trivial classes [{\leanfont Sha D e}]}
\nt{$\resInf{\mathcal D}{h}{\xi}$}{the restriction of the class~$\xi$ (Definition~\ref{d:classes}) [{\leanfont resInf D e ξ}]}
\nt{$\qzG{F}{\Omega}{\ell}{h}$}{$\QZl{\ell}$ with the action of $\GalG{F}{\Omega}$ through \m{\cyc{F}{\Omega}{\ell}{\cdot}^{-h}} [{\leanfont qzG F Ω l e}]}
\nt{$\EqHom{N}{C}{\pi}$}{the equivariant characters of Definition~\ref{d:eqhom} [{\leanfont EqHom N U ρ}]}
\nt{$b$}{a nonzero element of a $\Z_\ell$-module (Lemma~\ref{l:sep})}
\nt{$\Phi$, $T$}{the isomorphism of Theorem~A; in the summary, a twisted action}
\nt{$L$, $\eta$}{the base field and the generator in the first note}
\end{multicols}
\endgroup

% ------------------------------------------------------------------------------------- references
\head{References}
\begingroup\fontsize{9.4}{12.6}\selectfont\raggedright\renewcommand{\section}[2]{}
\begin{thebibliography}{NSW}\setlength{\itemsep}{2pt}
\bibitem[Al]{Al} L.\,Alpöge, announcement of proofs of the Riemann hypothesis and of the generalized Riemann
  hypothesis for Dirichlet characters, 20 November 2026, with a Lean archive;
  \href{\PinUrl}{public mirror of the archive at commit {\leanfont\Pin}}.
\bibitem[N1]{N1} \emph{Twisted cokernels in a cyclotomic tower}, Deformalization~1, riemann-rev,
  \href{https://nasqret.github.io/riemann-rev/digests/tower-bound.pdf}{\texttt{digests/tower-bound.pdf}}.
\bibitem[N2]{N2} \emph{Local nilpotence over a cyclotomic tower}, Deformalization~2, riemann-rev,
  \href{https://nasqret.github.io/riemann-rev/digests/local-nilpotence.pdf}{\texttt{digests/local-nilpotence.pdf}}.
\bibitem[NSW]{NSW} J.\,Neukirch, A.\,Schmidt and K.\,Wingberg, \emph{Cohomology of Number Fields}, 2nd ed.,
  Grundlehren der mathematischen Wissenschaften 323, Springer, 2008.
\bibitem[Se]{Se} J.-P.\,Serre, \emph{A Course in Arithmetic}, Graduate Texts in Mathematics 7, Springer, 1973.
\bibitem[Wa]{Wa} L.\,C.\,Washington, \emph{Introduction to Cyclotomic Fields}, 2nd ed., Graduate Texts in
  Mathematics 83, Springer, 1997.
\end{thebibliography}
\endgroup

\end{document}
